\chapter{Model calculations in dimension one and in the toric case}
\index{metric graph}\index{toric geometry!model calculation}
\label{chap:model-calculations}

This chapter turns the objects of
\crefrange{chap:valuations}{chap:topological-tools} into computations.  Curves
reduce model metrics to functions on metric graphs, while toric varieties
reduce invariant metrics to convex functions on a real vector space.  These
two settings will be the test suite for every later sign and normalization.

We retain the global convention $v=-\log|\cdot|$.  When $K$ is discretely
valued, the running uniformizer is normalized by $v(\pi)=1$.  Monge--Amp\`ere
measures are not divided by their total mass: a metric on a degree-$d$ line
bundle on a curve produces a measure of mass $d$.

\subsection*{What this chapter supplies to the rest of the book}

The chapter is deliberately the most concrete in \cref{part:foundations}, and
almost every statement in it is a computation carried to explicit numbers.
Four packages of results are assembled here and imported later.

\begin{enumerate}
\item \emph{The one-dimensional calculus.}  \Cref{sec:metric-graphs}
constructs the Laplacian of a piecewise affine function on a compact metric
graph, proves Green's identity, and solves the Poisson and Dirichlet
problems; \cref{sec:ch08-slope-formula} proves the slope formula relating the
Laplacian on a skeleton to the retraction of a principal divisor.
\Cref{part:cld} recovers the same
operator as $\dprimeddprime$ on the analytic curve
(\cref{prop:skeletal-curvature-curve}), \cref{part:bfj} uses it to solve the
Monge--Amp\`ere equation on curves (\cref{prop:graph-ma-solution}), and
\cref{part:bj} uses its version on the trivially valued tree
(\eqref{eq:bj-curve-laplacian}).

\item \emph{The projective line.}  \Cref{sec:berkovich-p1} makes the tree
picture of \cref{sec:berk-line} carry potentials: model metrics on
$\cO(d)$, their piecewise affine functions, their Chambert--Loir measures,
and the mass check against the degree.  It is the smallest place in which
the two computations of a curvature measure---by outgoing slopes and by
intersection numbers on a model---can be compared in full.

\item \emph{The trivially valued cone.}  \Cref{sec:trivial-valuation-cone}
computes the cone of semivaluations of a curve and of the affine plane,
its scaling action, its divisorial points, and the homogeneous potentials
attached to ideals.  \Cref{part:bj} works on exactly these objects; the
definitions there (\cref{def:trivial-analytification,def:homogeneous-function})
are the abstractions of what is computed here.

\item \emph{The toric dictionary.}  \Cref{sec:toric-convex} identifies
invariant metrics with convex functions on $N_{\mathbf R}$ and the
Monge--Amp\`ere measure with the real Alexandrov measure, up to the factor
$n!$ forced by \eqref{eq:appJ-toric-degree}.  Everything is carried out in
coordinates for $\mathbf P^1\times\mathbf P^1$, a Hirzebruch surface, and a
weighted projective plane.  \Cref{thm:toric-bfj-alexandrov,thm:bj-toric-psh-convex}
are the two later avatars.
\end{enumerate}

A reader who has done the computations below by hand can check any later sign
without re-deriving it: the four normalization tests collected at the end of
\cref{sec:model-translation-table} determine every convention in the book.

\subsection*{Dependencies}

\Cref{sec:metric-graphs} is self-contained apart from the model theory of
\cref{chap:models} and the skeleton theory of \cref{chap:skeleta}, both used
only in \cref{sec:ch08-slope-formula}.  \Cref{sec:berkovich-p1} uses the
classification of points of \cref{thm:berk-a1-classification} and the model
measures of \cref{sec:metrics-cl-measure}.  \Cref{sec:trivial-valuation-cone}
uses only \cref{sec:berk-trivial}.  \Cref{sec:toric-convex} uses convexity
in the form recorded in \cref{app:convexity} and the model measures again.
The three running examples of \cref{sec:ch01-running-examples} are recomputed
throughout, and are never redefined here: the nodal degeneration is
\cref{sec:ch01-running-nodal}, the toric line is
\cref{sec:ch01-running-toric}, and the product test configuration is
\cref{sec:ch01-running-tc}.

\section{Metric graphs and their Laplacians}
\label{sec:metric-graphs}

A skeleton of a curve is a finite metric graph, and every model metric on a
line bundle over that curve becomes a piecewise affine function on it.  This
section develops the potential theory of such functions from scratch.
Nothing in it uses algebraic geometry until
\cref{sec:ch08-slope-formula}.

\subsection{Graphs, models, and piecewise affine functions}
\label{sec:ch08-graph-basics}

\begin{definition}
A \emph{compact metric graph} $\Gamma$ is a finite connected graph whose
edges are assigned positive lengths and realized as closed intervals, with
the induced path metric.  At $p\in\Gamma$, the finite set of germs of
nonconstant geodesic segments issuing from $p$ is denoted $T_p\Gamma$ and its
elements are called tangent directions.
\end{definition}

The data of a finite graph with edge lengths is called a \emph{model} of the
topological space $\Gamma$; the space itself, with its metric, does not
remember which points were declared to be vertices.  Inserting a new vertex
in the interior of an edge of length $\ell$, at distance $\ell'$ from one of
its endpoints, replaces that edge by two edges of lengths $\ell'$ and
$\ell-\ell'$ and changes neither $\Gamma$ nor its path metric; such an
operation is a \emph{subdivision}.  Two models of the same metric graph
always admit a common subdivision, obtained by taking the union of the two
vertex sets.  This trivial remark is the combinatorial shadow of a
substantial geometric fact: distinct models of a curve may have the same
skeleton, and models related by vertical blowup subdivide it
(\cref{ex:snc-chain-models,thm:skeleton-subdivision}).

\begin{example}[Valence and tangent directions]
\label{ex:ch08-valence}
If $p$ is a vertex of a model of $\Gamma$, then $\#T_p\Gamma$ is its valence
in that model.  If $p$ lies in the interior of an edge, then
$\#T_p\Gamma=2$, whatever model is used, and $p$ is a two-valent vertex of
any subdivision containing it.  An endpoint of a segment has
$\#T_p\Gamma=1$.  Thus for $\Gamma=[0,\ell]$ one has $\#T_0\Gamma=1$,
$\#T_\ell\Gamma=1$, and $\#T_p\Gamma=2$ for $0<p<\ell$; for the circle
$\mathbf R/\ell\mathbf Z$ every point has exactly two tangent directions.
\end{example}

A continuous function $f$ on $\Gamma$ is \emph{piecewise affine} if, after a
finite subdivision, its restriction to each edge is affine.  For
$\nu\in T_p\Gamma$, let $d_\nu f(p)$ be the outgoing slope of $f$ in the
direction $\nu$.

Explicitly, if $\gamma:[0,\varepsilon)\to\Gamma$ is the unit-speed geodesic
representing $\nu$, then
\begin{equation}
\label{eq:ch08-outgoing-slope}
 d_\nu f(p)=\lim_{s\downarrow0}\frac{f(\gamma(s))-f(p)}{s},
\end{equation}
which exists for piecewise affine $f$ because $f\circ\gamma$ is affine near
$0$.  We write $\PL(\Gamma)$ for the real vector space of piecewise affine
functions on $\Gamma$; it contains the constants and is stable under
$\max$, $\min$, and sums.

\begin{definition}
The graph Laplacian of a piecewise affine function is the signed atomic
measure
\begin{equation}
\label{eq:graph-laplacian}
             \Delta_\Gamma f
             =\sum_{p\in\Gamma}
                \left(\sum_{\nu\in T_p\Gamma}d_\nu f(p)\right)\delta_p.
\end{equation}
Only vertices and slope-break points contribute.
\end{definition}

Our sign convention makes a convex corner positive.  For example,
$f(t)=|t|$ on a neighborhood of $0$ in a line has
$\Delta f=2\delta_0$.

Formula \eqref{eq:graph-laplacian} is the operator $\ddc$ of
\cite[\S1.2.1]{Thu05} on an integral polyhedron of dimension at most one,
specialized to the constant weight $1$.  Thuillier writes
$\ddc f=\sum_{x}\lambda_x(f)\delta_x$ with
$\lambda_x(f)=\sum_{t\in\mathrm T_xS}w_x(t)\lambda_{x,t}(f)$, the inner sum
running over the tangent directions at $x$ and $\lambda_{x,t}(f)$ being the
outgoing slope of $f$ along $t$; the weights $w_x(t)$ are the residual degrees
of the nodes when the polyhedron is the reduction graph of a nodal curve
\cite[Rem.~1.2.4]{Thu05}, and they are all equal to $1$ in the situation of
\cref{sec:ch08-slope-formula}, where $\widetilde K$ is algebraically closed.
The whole of \cref{sec:metric-graphs} is his \cite[\S1.2]{Thu05}, and we
quote his numbering statement by statement below.

Two points must be checked before \eqref{eq:graph-laplacian} can be used:
that the sum is finite, and that it does not depend on the model chosen to
compute it.

\begin{lemma}
\label{lem:ch08-laplacian-well-defined}
Let $f\in\PL(\Gamma)$.
\begin{enumerate}
\item The set of $p\in\Gamma$ with
$\sum_{\nu\in T_p\Gamma}d_\nu f(p)\neq0$ is finite, so
\eqref{eq:graph-laplacian} defines a signed atomic measure of finite total
variation.
\item The measure $\Delta_\Gamma f$ does not depend on any choice of model:
it is computed from $f$ and the metric alone.
\item $\Delta_\Gamma$ is linear, and $\Delta_\Gamma f=0$ if and only if $f$
is constant.
\end{enumerate}
\end{lemma}

\begin{proof}
(1)  Fix a model in which $f$ is affine on every edge.  If $p$ lies in the
interior of an edge $e$ on which $f$ is affine with slope $c$ (computed in
some orientation of $e$), then $T_p\Gamma$ consists of the two directions
along $e$, with outgoing slopes $c$ and $-c$; their sum vanishes.  Hence the
support of $\Delta_\Gamma f$ is contained in the finite vertex set.

(2)  The quantity $\sum_{\nu\in T_p\Gamma}d_\nu f(p)$ is defined by
\eqref{eq:ch08-outgoing-slope} purely in terms of the metric space
$\Gamma$, its tangent directions at $p$, and the germ of $f$ at $p$.  No
model enters.  Part~(1) shows only that a model exhibits the finiteness; by
the previous paragraph, adding vertices adds only points of zero mass.

(3)  Linearity is clear from \eqref{eq:ch08-outgoing-slope}.  If $f$ is
constant then every outgoing slope vanishes.  Conversely suppose
$\Delta_\Gamma f=0$ and let $p_0\in\Gamma$ be a point at which the continuous
function $f$ attains its maximum; we may take $p_0$ to be a vertex of a model
in which $f$ is edgewise affine, since the maximum of an affine function on an
edge is attained at an endpoint.  All outgoing slopes at $p_0$ are
$\leq0$, and their sum is $0$, so all of them vanish; therefore $f$ is
constant on every edge issuing from $p_0$, and every neighbouring vertex is
again a maximum point.  Since $\Gamma$ is connected, $f$ is constant.
\end{proof}

The last argument is the graph maximum principle, isolated below in the form
used repeatedly in \cref{part:bfj}.

\begin{proposition}
\label{prop:ch08-graph-maximum}
Let $f\in\PL(\Gamma)$ and let $S\subset\Gamma$ be a nonempty finite set
containing the support of the negative part of $\Delta_\Gamma f$.  Then
\[
                 \max_\Gamma f=\max_{S}f .
\]
Consequently, if $\Delta_\Gamma f\geq0$ on all of $\Gamma$, then $f$ is
constant.
\end{proposition}

\begin{proof}
Put $M=\max_\Gamma f$ and $U=\{f=M\}$, a nonempty closed subset.  Let
$p_0\in U\setminus S$.  Choose a model in which $f$ is affine on every edge
and $S\cup\{p_0\}$ lies in the vertex set.  All outgoing slopes of $f$ at
$p_0$ are $\leq0$ because $f\leq M=f(p_0)$, while their sum is
$(\Delta_\Gamma f)(\{p_0\})\geq0$ because $p_0\notin S$.  Hence every
outgoing slope at $p_0$ vanishes, so $f\equiv M$ on every edge incident to
$p_0$; that is, $U$ contains the closed star of $p_0$, which is a
neighbourhood of $p_0$.

Suppose now that $U\cap S=\varnothing$.  The previous paragraph applies at
every point of $U$, so $U$ is open as well as closed and nonempty; since
$\Gamma$ is connected, $U=\Gamma$ and $f\equiv M$.  But then
$U\supseteq S\neq\varnothing$, contradicting $U\cap S=\varnothing$.
Therefore $U\cap S\neq\varnothing$, which is the assertion.

For the last statement, $\Delta_\Gamma f\geq0$ together with
$\Delta_\Gamma f(\Gamma)=0$, proved as \eqref{eq:lap-mass-zero} below, forces
$\Delta_\Gamma f=0$, and \cref{lem:ch08-laplacian-well-defined}\,(3) applies.
\end{proof}

In the harmonic case $\Delta_\Gamma f=0$ off $S$ this is
\cite[Prop.~1.2.9]{Thu05}, in the global form \cite[Cor.~1.2.10]{Thu05}; the
statement above is its subharmonic refinement, obtained by his argument
verbatim.  The consequence that a harmonic function on a compact connected
polyhedron with at most one boundary point is constant is
\cite[Cor.~1.2.11]{Thu05}, and it is the polyhedral shadow of the fact that a
psh function on a disc is determined by its value at the boundary point of
the disc.

\begin{proposition}\label{prop:graph-green}
For piecewise affine $f,g$ on a compact metric graph,
\begin{align}
 \Delta_\Gamma f(\Gamma)&=0,\label{eq:lap-mass-zero}\\
 \int_\Gamma f\,\rd \Delta_\Gamma g
 &=\int_\Gamma g\,\rd \Delta_\Gamma f
 =-\int_\Gamma f'g'\,\rd s.\label{eq:graph-green}
\end{align}
Here the last integral is the sum of the ordinary edge integrals and is
independent of the chosen subdivision.
\end{proposition}

\begin{proof}
On every oriented edge $e=[a,b]$, the outgoing slope at $a$ is $f'_e$ and
the outgoing slope at $b$ is $-f'_e$; these contributions cancel in the
total mass, including after subdivision.  Integration by parts on each
affine piece gives
\[
 -\int_e f'g'\,\rd s
 =f(a)d_eg(a)+f(b)d_{\bar e}g(b),
\]
where $e$ and $\bar e$ are the two outgoing directions.  Summing over all
pieces yields the first equality in \eqref{eq:graph-green}; symmetry gives
the second.
\end{proof}

The symmetry of the pairing, that is the first equality in
\eqref{eq:graph-green}, is \cite[Prop.~1.2.13]{Thu05}, and the mass identity
\eqref{eq:lap-mass-zero} is his \cite[Cor.~1.2.14]{Thu05}.  Both are used in
exactly the same way there as here: symmetry is what allows $\ddc$ to be
transposed to an operator on distributions, and vanishing of the total mass
is what forces a curvature measure to have the mass of its line bundle.

\begin{example}[The segment]
Let $\Gamma=[0,\ell]$ and $f(t)=t$.  The outgoing slopes give
\[
                         \Delta_\Gamma f=\delta_0-\delta_\ell.
\]
Consequently
$\int f\,\rd \Delta f=-\ell=-\int_0^\ell|f'|^2\,\rd t$, confirming both the endpoint
sign and the Green identity.
\end{example}

\begin{corollary}
\label{cor:ch08-dirichlet-energy}
For $f\in\PL(\Gamma)$ set
\[
                 D(f)=\int_\Gamma|f'|^2\,\rd s\ \geq 0 .
\]
Then $D(f)=-\int_\Gamma f\,\rd\Delta_\Gamma f$, and $D(f)=0$ if and only if
$f$ is constant.  In particular the symmetric bilinear form
$(f,g)\mapsto-\int_\Gamma f\,\rd\Delta_\Gamma g$ is positive semidefinite
with kernel the constants.
\end{corollary}

\begin{proof}
Take $g=f$ in \eqref{eq:graph-green}.  The integrand $|f'|^2$ is nonnegative
and vanishes identically exactly when every edge slope vanishes, that is,
when $f$ is constant.  Symmetry and bilinearity are
\eqref{eq:graph-green} again.
\end{proof}

The sign in \cref{cor:ch08-dirichlet-energy} is worth memorizing.  The
operator $\Delta_\Gamma$ is \emph{negative} semidefinite, exactly like the
usual Laplacian $\partial^2/\partial t^2$ on an interval and unlike the
positive combinatorial graph Laplacian $L=D-A$ of spectral graph theory.
Our normalization is chosen so that a convex corner carries positive mass,
because a convex corner of a metric weight is where positive curvature sits.

\begin{remark}[Thuillier's normalization, and the factor that separates it
from Rumely's]
\label{rem:ch08-thuillier-normalization}
Three conventions have to be matched before \cite{Thu05} can be quoted, and
only the third produces a factor.

\emph{Sign.}  Thuillier's $\ddc$ on a one-dimensional integral polyhedron is
the sum of outgoing slopes, exactly as in \eqref{eq:graph-laplacian}, and his
subharmonic functions restrict to \emph{convex} functions on a skeleton
\cite[Rem.~3.1.7 and Prop.~3.4.6]{Thu05}.  A metric on a line bundle is
subharmonic for him when $-\log\|s\|$ is subharmonic for every invertible
local section \cite[Def.~3.4.17]{Thu05}, and its curvature is the unique
degree-one current restricting to $\ddc(-\log\|s\|)$ on the locus where such
an $s$ exists, hence $c_1=\ddc(-\log\|s\|)+\delta_{\divis(s)}$ for a
meromorphic section; that is exactly
\eqref{eq:appJ-cld-section-pl}, and no sign conversion is needed.  We use
the convention fixed by the cited definition and curvature computation,
without asserting an erratum in the source text.

\emph{Edge lengths.}  At a node with formal equation $x_1x_2=a$ the integral
structure of \cite[\S3.6.3]{Thu05} gives the corresponding edge the length
$-\log|a|$.  Under the convention $v=-\log|\cdot|$ of
\cref{app:normalizations} this is $v(a)$, which is the length used in
\cref{sec:ch08-slope-formula}; with $v(\pi)=1$ the two agree on the nose and
no rescaling occurs.

\emph{The one genuine factor.}  Rumely's reduction graph of a well adjusted
model over a finite extension $K/\mathbf Q_p$ gives each edge the length
$1/e(K/\mathbf Q_p)$ rather than $-\log|a|=e(K/\mathbf Q_p)^{-1}\log p$.
Consequently Thuillier's operator and Rumely's differ by a constant,
\[
   \ddc=(\log p)^{-1}\,\ddc_{\mathrm R},
\]
and, the Green function of a metric graph scaling linearly in the edge
lengths, his smooth Green functions are $\log p$ times Rumely's; this is
\cite[\S5.1]{Thu05}.  A reader who takes edge lengths from the
$p$-adic-normalized literature, where $|p|=p^{-1}$, and Laplacians from this
book must therefore insert $\log p$ once.  We never do: the absolute value is
normalized by $v(\pi)=1$ throughout, so $\log p$ never appears.
\end{remark}

\subsection{Harmonic functions and the Dirichlet problem}
\label{sec:ch08-harmonic}
\index{harmonic function!on a metric graph}

\begin{definition}
Let $S\subset\Gamma$ be a finite set.  A function $f\in\PL(\Gamma)$ is
\emph{harmonic outside $S$} if $\Delta_\Gamma f$ is supported on $S$, and
\emph{harmonic} if $\Delta_\Gamma f=0$.
\end{definition}

By \cref{lem:ch08-laplacian-well-defined}\,(3), the only harmonic functions on
a compact metric graph are the constants.  The interesting statement is
therefore the boundary-value problem, in which $S$ plays the role of the
boundary.

\begin{proposition}
\label{prop:ch08-dirichlet}
Let $S=\{p_1,\ldots,p_N\}\subset\Gamma$ be a finite nonempty set and let
$c_1,\ldots,c_N\in\mathbf R$.  There is exactly one $f\in\PL(\Gamma)$ that is
harmonic outside $S$ and satisfies $f(p_i)=c_i$ for all $i$.  It satisfies
\[
   \min_i c_i\ \leq\ f\ \leq\ \max_i c_i \qquad\text{on }\Gamma,
\]
and it minimizes the Dirichlet energy $D$ among all piecewise affine
functions with the prescribed values on $S$.
\end{proposition}

\begin{proof}
Fix a model whose vertex set $V$ contains $S$, and look for $f$ affine on
every edge.  Such an $f$ is determined by its values on $V$, and for a vertex
$p\in V\setminus S$ the harmonicity condition reads
\begin{equation}
\label{eq:ch08-discrete-harmonic}
 \sum_{q\sim p}\frac{f(q)-f(p)}{\ell_{pq}}=0 ,
\end{equation}
the sum being over the edges at $p$, with $\ell_{pq}$ the length of the edge
and multiple edges counted with multiplicity.  This is a linear system of
$\#(V\setminus S)$ equations in the same number of unknowns
$(f(p))_{p\in V\setminus S}$.  For uniqueness it suffices to prove that the
homogeneous system has only the zero solution: if $f|_S=0$ and $f$ is
harmonic outside $S$, then \cref{prop:ch08-graph-maximum} applied to $f$ and
to $-f$ gives $\max_\Gamma f=\max_Sf=0$ and $\min_\Gamma f=0$, so $f=0$.  A
square linear system with trivial kernel is invertible, which gives existence
as well.  Passing to a subdivision does not change the solution, because a
two-valent vertex with vanishing outgoing-slope sum is exactly a point at
which the two adjacent affine pieces have equal slopes.

The bound $\min_ic_i\leq f\leq\max_ic_i$ is
\cref{prop:ch08-graph-maximum} applied to $\pm f$.  For the variational
characterization, let $g\in\PL(\Gamma)$ agree with $f$ on $S$ and set
$h=g-f$, so $h|_S=0$.  Then
\[
 D(g)=D(f)+2\int_\Gamma f'h'\,\rd s+D(h)
     =D(f)-2\int_\Gamma h\,\rd\Delta_\Gamma f+D(h)
\]
by \eqref{eq:graph-green}.  The middle term vanishes because
$\Delta_\Gamma f$ is supported on $S$ and $h|_S=0$.  Hence
$D(g)=D(f)+D(h)\geq D(f)$, with equality if and only if $h$ is constant,
that is, $h\equiv0$.
\end{proof}

The bijectivity assertion in \cref{prop:ch08-dirichlet}---the restriction map
from the space of functions harmonic outside $S$ to $\mathbf R^S$ is an
isomorphism---is \cite[Prop.~1.2.15]{Thu05}, and the resulting local
solvability of the Dirichlet problem on an integral polyhedron is
\cite[Rem.~1.2.16]{Thu05}.  That remark is the combinatorial input from which
Thuillier deduces the analytic statement: every point of a smooth strictly
$K$-analytic curve has a fundamental system of neighbourhoods on which the
Dirichlet problem for continuous boundary data is uniquely solvable
\cite[\S3.1.1 and Cor.~3.1.21]{Thu05}.  With the maximum principle
\cite[Prop.~3.1.1]{Thu05} and the Harnack principle
\cite[Prop.~3.1.2]{Thu05}, this is precisely the axiomatic package of Bauer,
Brelot and Doob, and Thuillier transcribes Perron's method into it
\cite[\S3.1.3]{Thu05}, the barriers being supplied at every boundary point of
type~$\mathrm{II}$, $\mathrm{III}$ or $\mathrm{IV}$ by
\cite[Prop.~3.1.20]{Thu05}.  Nothing beyond the finite-dimensional linear
algebra above is needed on a metric graph.

\begin{example}[Two boundary points on a segment]
\label{ex:ch08-segment-dirichlet}
Take $\Gamma=[0,\ell]$, $S=\{0,\ell\}$, and prescribed values $A$ and $B$.
The unique solution is the affine function
$f(t)=A+(B-A)t/\ell$, whose Laplacian is
\[
 \Delta_\Gamma f=\frac{B-A}{\ell}\delta_0+\frac{A-B}{\ell}\delta_\ell ,
\]
and whose Dirichlet energy is $D(f)=(B-A)^2/\ell$.  This is the computation
underlying \cref{ex:edge-mass-transport} below: a long edge transports mass
cheaply, a short edge expensively.
\end{example}

Let $\mu_0$ be a positive atomic measure on $\Gamma$ of mass $d$.  A graph
potential $\varphi$ changes it to
\begin{equation}
\label{eq:graph-curvature}
                         \mu_\varphi=\mu_0+\Delta_\Gamma\varphi.
\end{equation}
By \eqref{eq:lap-mass-zero}, the mass remains $d$.  In the model-theoretic
interpretation, $\mu_0$ is the curvature measure of a reference model metric
and $\varphi$ is a model function.  Semipositivity is exactly the condition
$\mu_\varphi\geq0$.

\begin{example}[Moving mass across one edge]
\label{ex:edge-mass-transport}
Let $\mu_0=d_0\delta_0+d_1\delta_\ell$ with $d_0,d_1\geq0$.  Prescribe
$\varphi(0)=A$ and $\varphi(\ell)=B$, and make $\varphi$ affine.  Then
\[
 \Delta\varphi=
 \frac{B-A}{\ell}\delta_0+
 \frac{A-B}{\ell}\delta_\ell,
\]
and hence
\[
 \mu_\varphi=
 \left(d_0+\frac{B-A}{\ell}\right)\delta_0+
 \left(d_1+\frac{A-B}{\ell}\right)\delta_\ell.
\]
The perturbed metric is semipositive if and only if
\[
                         -d_0\ell\leq B-A\leq d_1\ell.
\]
Thus nefness on the two components becomes a pair of slope inequalities.
\end{example}

\subsection{Divisors and tropical rational functions}
\label{sec:ch08-tropical-rational}

An integral piecewise affine function has integral slopes.  Its Laplacian is
then a divisor of degree zero on $\Gamma$.  The equality
\[
                         \operatorname{div}_\Gamma(f)=\Delta_\Gamma f
\]
is the graph-theoretic form of the slope formula.  Constants are the kernel
of $\Delta_\Gamma$ because $\Gamma$ is connected.

\begin{proposition}\label{prop:graph-poisson}
For every finitely supported signed atomic measure
$\nu=\sum_pa_p\delta_p$ on a compact metric graph with total mass zero,
there is a continuous piecewise affine function $f$, unique up to an additive
constant, such that $\Delta_\Gamma f=\nu$.
\end{proposition}

\begin{proof}
Choose a vertex set containing the support of $\nu$.  On a vertex $p$, the
outgoing-slope operator is
\[
 f\longmapsto\sum_{q\sim p}\frac{f(q)-f(p)}{\ell_{pq}},
\]
the negative of the usual weighted combinatorial Laplacian.  Its image is
the hyperplane of vectors with coordinate sum zero, and its kernel consists
of constants.  Solve the resulting finite linear system and extend the
solution affinely over each edge.
\end{proof}

Restated, \cref{prop:graph-poisson} together with \eqref{eq:lap-mass-zero}
and \cref{lem:ch08-laplacian-well-defined}\,(3) says that
\[
 0\longrightarrow\mathbf R\longrightarrow
 \{f\in\PL(\Gamma):\Supp\Delta_\Gamma f\subseteq S\}
 \xrightarrow{\ \Delta_\Gamma\ }
 \mathrm M(S)\xrightarrow{\ \int_\Gamma\ }\mathbf R\longrightarrow0
\]
is exact, $\mathrm M(S)$ denoting the signed measures supported on the finite
set $S$; this is \cite[Prop.~1.2.17]{Thu05}.  The elementary solutions
$h_{p_0,p}$ normalized by $h_{p_0,p}(p_0)=0$ and
$\Delta_\Gamma h_{p_0,p}=\delta_p-\delta_{p_0}$ are
\cite[Lem.~1.2.18]{Thu05}; they are the Green functions written down in
closed form in \cref{sec:ch08-green-functions}.

\subsection{Green functions, computed}
\label{sec:ch08-green-functions}
\index{Green function!on a metric graph}

\Cref{prop:graph-poisson} is an existence statement; the two cases used
throughout the book can be written down in closed form.  Both computations
below are carried out entirely by counting outgoing slopes.

\begin{example}[Green function of a segment]
\label{ex:ch08-green-segment}
Let $\Gamma=[0,\ell]$ and fix $q\in(0,\ell)$.  Look for the piecewise affine
$g$ with $g(0)=0$, affine on $[0,q]$ and on $[q,\ell]$, and with
\[
                    \Delta_\Gamma g=\alpha\delta_0+\beta\delta_\ell-\delta_q .
\]
Writing $g(q)=c$ and $g(\ell)=b$, the three outgoing-slope sums are
\[
 \alpha=\frac{c}{q},\qquad
 \beta=\frac{c-b}{\ell-q},\qquad
 -1=\frac{-c}{q}+\frac{b-c}{\ell-q} .
\]
The third equation is the sum of the first two with a sign, so
$\alpha+\beta=1$, as \eqref{eq:lap-mass-zero} requires.  Choosing the
normalization $b=0$, so that $g$ vanishes at both endpoints, gives
$c=q(\ell-q)/\ell$ and
\begin{equation}
\label{eq:ch08-green-segment}
 g(s)=
 \begin{cases}
  \dfrac{(\ell-q)s}{\ell}, & 0\leq s\leq q,\\[2mm]
  \dfrac{q(\ell-s)}{\ell}, & q\leq s\leq \ell,
 \end{cases}
 \qquad
 \Delta_\Gamma g=\frac{\ell-q}{\ell}\delta_0+\frac{q}{\ell}\delta_\ell
                 -\delta_q .
\end{equation}
The endpoint weights $(\ell-q)/\ell$ and $q/\ell$ are barycentric: a unit
charge at $q$ is balanced by the two endpoints in inverse proportion to
their distances.  The Dirichlet energy is
\[
 D(g)=q\left(\frac{\ell-q}{\ell}\right)^2
      +(\ell-q)\left(\frac{q}{\ell}\right)^2
     =\frac{q(\ell-q)}{\ell},
\]
which agrees with $-\int g\,\rd\Delta_\Gamma g=g(q)=q(\ell-q)/\ell$, as
\cref{cor:ch08-dirichlet-energy} predicts.  For $\ell=1$ this is
\eqref{eq:segment-green}, with maximal energy $\tfrac14$ at $q=\tfrac12$.
\end{example}

\begin{example}[Green function of a circle]
\label{ex:ch08-green-circle}
\index{circle!metric graph}
Let $\Gamma=\mathbf R/\ell\mathbf Z$ be the circle of circumference $\ell$,
and let $0<c<\ell$.  We solve
\[
                 \Delta_\Gamma g=\delta_0-\delta_c
\]
with $g$ affine on the two arcs $[0,c]$ and $[c,\ell]$.  Normalize $g(0)=0$
and put $g(c)=\kappa$.  At the point $0$ the two tangent directions run into
the two arcs, of lengths $c$ and $\ell-c$, and both reach the value $\kappa$;
so
\[
 (\Delta_\Gamma g)(\{0\})
   =\frac{\kappa}{c}+\frac{\kappa}{\ell-c}
   =\kappa\,\frac{\ell}{c(\ell-c)} .
\]
Setting this equal to $1$ gives $\kappa=c(\ell-c)/\ell$, and by symmetry the
mass at $c$ is $-1$.  Hence
\begin{equation}
\label{eq:ch08-green-circle}
 g(s)=
 \begin{cases}
  \dfrac{(\ell-c)s}{\ell},& 0\leq s\leq c,\\[2mm]
  \dfrac{c(\ell-s)}{\ell},& c\leq s\leq\ell,
 \end{cases}
 \qquad
 g(c)=\frac{c(\ell-c)}{\ell},
 \qquad
 D(g)=\frac{c(\ell-c)}{\ell} .
\end{equation}
The formula is literally the same as \eqref{eq:ch08-green-segment}, but its
meaning is different: on the circle there are no endpoints to absorb charge,
so the balancing charge must be a genuine second point mass.  For $\ell=1$
and $c=\tfrac12$ one gets $g(\tfrac12)=\tfrac14$ and $D(g)=\tfrac14$.  The
circle is the skeleton of a degenerating elliptic curve
(\cref{ex:ch08-tate-curve}), and \eqref{eq:ch08-green-circle} is the
non-Archimedean N\'eron function of the divisor $(0)-(c)$ on it.
\end{example}

\begin{remark}[Why the answers are rational]
Both Green functions have values in $\mathbf Q$ whenever the edge lengths and
the charge positions are rational.  This is a general feature of
\cref{prop:graph-poisson}: the linear system \eqref{eq:ch08-discrete-harmonic}
has rational coefficients.  It is the graph reflection of the fact that
divisorial points, which have rational coordinates on a skeleton, suffice for
all model computations (\cref{thm:skeleton-pl-realization}).
\end{remark}

\section{Skeleta of curves, the slope formula, and the retraction}
\label{sec:ch08-slope-formula}
\index{slope formula}\index{retraction!and Laplacian}

We now connect \cref{sec:metric-graphs} to geometry.  Throughout this
section, $K$ is complete and discretely valued with $v(\pi)=1$ and
algebraically closed residue field $\widetilde K$, $R=K^\circ$, and $C$ is a
smooth, connected, projective curve over $K$ admitting a regular semistable
model.  Recall from \cref{chap:models,chap:skeleta} what this means.

\begin{definition}[Semistable model of a curve]
\label{def:ch08-semistable-curve-model}
A \emph{regular semistable model} of $C$ is a regular, flat, projective
$R$-scheme $\cX$ with generic fibre $C$ whose special fibre
$\cX_s=\sum_YY$ is reduced---so every multiplicity is $b_Y=1$---with smooth
irreducible components meeting transversally in ordinary double points, at
which the completed local ring of $\cX$ is $R[[x,y]]/(xy-\pi)$.
\end{definition}

Such a model is in particular an SNC model in the sense of
\cref{chap:skeleta}, and its dual complex $\Delta_{\cX}$ is a compact metric
graph: one vertex $v_Y$ for each component $Y$, one edge of length one for
each node, and $\Sk(\cX)=\iota_{\cX}(\Delta_\cX)\subset C^{\An}$ by
\cref{thm:skeleton-embedding}.  We write $\Gamma=\Sk(\cX)$ and
$r=r_{\cX}:C^{\An}\to\Gamma$ for the retraction of
\cref{prop:skeleton-retraction}.

Concretely, at a node with formal equation $xy=\pi$ the corresponding edge is
$\{v_s:0\leq s\leq1\}$ with $v_s(x)=s$ and $v_s(y)=1-s$
(\cref{sec:ch01-running-nodal}), and $\iota_{\cX}$ identifies it with the
unit interval; the retraction sends a point $z$ of the corresponding open
annulus to $v_{v_z(x)}$.  For a node with equation $xy=a$, $a\in R$, the
edge has length $v(a)$.

\subsection{The retraction is a maximum}

\begin{proposition}
\label{prop:ch08-retraction-max}
Let $\cX$ be a regular semistable model of $C$ and let
$z\in C^{\An}$.  For every $f\in\cO_{\cX,\xi}$, where $\xi$ is the centre of
$z$ on $\cX$,
\[
                  |f(z)|\ \leq\ |f(r(z))| ,
                  \qquad\text{equivalently}\qquad
                  v_z(f)\ \geq\ v_{r(z)}(f) .
\]
Consequently, for a vertical fractional ideal $\mathfrak a$ on $\cX$,
$\log|\mathfrak a|(z)\leq\log|\mathfrak a|(r(z))$, and a model function
$\varphi_D$ determined on $\cX$ satisfies $\varphi_D\circ r=\varphi_D$.
\end{proposition}

\begin{proof}
If $\xi$ is the generic point of a component $Y$, then $r(z)=v_Y$ is the
divisorial point of $Y$, $\cO_{\cX,\xi}$ is a discrete valuation ring with
uniformizer $\pi$, and $v_{r(z)}(f)=\ord_Y(f)$ is the least $m$ with
$f\in\pi^m\cO_{\cX,\xi}$.  Since $v_z$ is a semivaluation with
$v_z(\pi)=1$ and $v_z\geq0$ on $\cO_{\cX,\xi}$, writing $f=\pi^mu$ with $u$ a
unit gives $v_z(f)=m+v_z(u)\geq m=v_{r(z)}(f)$.

If $\xi$ is a node, then $\cX$ is regular of dimension two at $\xi$ and, by
\cref{def:ch08-semistable-curve-model},
$\widehat{\cO}_{\cX,\xi}\simeq R[[x,y]]/(xy-\pi)$, in which every element is
uniquely
\[
 f=\sum_{i\geq0}a_ix^i+\sum_{j\geq1}b_jy^j,
 \qquad a_i,b_j\in R .
\]
In equal characteristic zero a Cohen coefficient field turns this into
$\kappa(\xi)[[x,y]]$ with $\pi=u\,xy$, as in \cref{sec:skeleton-snc}; the
presentation over $R$ is the one available in every characteristic, and the
monomial valuation below is likewise characteristic-free, being built from
the $\alpha$-weighted filtration (\cref{sec:skeleton-monomial}).  Put
$\alpha=(v_z(x),v_z(y))$; then $\alpha_1+\alpha_2=v_z(\pi)=1$ and, by
\cref{prop:skeleton-retraction}, $r(z)$ is the monomial valuation $v_\alpha$,
whose value on $f$ is
\[
 v_{r(z)}(f)=\min\Bigl\{
   \min_{a_i\neq0}\bigl(v(a_i)+i\alpha_1\bigr),\
   \min_{b_j\neq0}\bigl(v(b_j)+j\alpha_2\bigr)\Bigr\} .
\]
Since $z$ has centre $\xi$, both coordinates of $\alpha$ are strictly
positive: if, say, $v_z(x)=0$, then $x$ would be a unit for $z$ and the
centre of $z$ would not lie on the component $\{x=0\}$.  Hence
$v_z(\mathfrak m_\xi^n)\geq n\min\{\alpha_1,\alpha_2\}\to+\infty$, so $v_z$
extends continuously to the $\mathfrak m_\xi$-adic completion.  Applying the
ultrametric inequality to the partial sums of the series and passing to the
limit gives $v_z(f)\geq v_{r(z)}(f)$.

Finally, if $\mathfrak a\subseteq\cO_{\cX}$ is generated near $\xi$ by
$f_1,\ldots,f_N$, then
$\log|\mathfrak a|(z)=-\min_\ell v_z(f_\ell)\leq-\min_\ell v_{r(z)}(f_\ell)
=\log|\mathfrak a|(r(z))$.  A vertical \emph{fractional} ideal is the unit
ideal on the generic fibre, so $\pi^N\mathfrak a\subseteq\cO_{\cX}$ for some
$N\geq0$ and $\log|\mathfrak a|=N+\log|\pi^N\mathfrak a|$; the inequality
therefore persists.  For a vertical Cartier divisor $D$ the argument is not
the inequality applied twice --- if $D$ is effective and $\xi\in\Supp D$ then
its local equation $f$ lies in $\cO_{\cX,\xi}$ but $f^{-1}$ does not --- but
a direct computation: by \Cref{lem:snc-normal-form}\,(1) one has
$f=u\prod_jz_j^{a_j}$ with $u$ a unit and the $z_j$ local equations of the
components of $\cX_s$ through $\xi$, so $v_z(u)=0$ and
$v_z(z_j)=\alpha_j=v_{r(z)}(z_j)$ by the construction of $r$; both sides
equal $\sum_ja_j\alpha_j$, and $\varphi_D\circ r=\varphi_D$.  This is also
the identity recorded in \cref{sec:skeleton-model-functions}.
\end{proof}

\begin{remark}
\Cref{prop:ch08-retraction-max} says that $\Gamma$ contains the Shilov
boundary of every formal fibre: the supremum of $|f|$ over the fibre
$r^{-1}(p)$ is attained at $p$ itself.  This is the reason that a model
metric, whose weight is constant on the fibres of $r$, loses no information
by being restricted to $\Gamma$, and the reason that Chambert--Loir measures
of model metrics are supported on $\Gamma$.
\end{remark}

\begin{remark}[The analytic form of \cref{prop:ch08-retraction-max}]
\label{rem:ch08-retraction-thuillier}
On the polyhedral side, the statement that a retraction
$\tau:S\to S'$ onto a polyhedral subdomain induces an isomorphism between the
harmonic functions on $S'$ and those on $S$ is \cite[Prop.~1.2.20]{Thu05};
this is what allows the entire potential theory of a curve to be transported
from its skeleta, and Thuillier's sheaf of harmonic germs is defined by
$\mathrm H(\mathcal X)=\tau_{\mathcal X}^{*}\,\mathrm
H\bigl(S(\mathcal X),\partial\mathcal X_\eta\bigr)$
\cite[Def.~2.3.1 and Def.~2.3.16]{Thu05}.  On the analytic side, the
inequality of \cref{prop:ch08-retraction-max} holds not only for
$\log|\mathfrak a|$ but for every subharmonic function: if $u$ is subharmonic
on $\mathcal X_\eta$ and $U=u|_{S(\mathcal X)}$, then $U$ is continuous and
convex and $u\leq U\circ\tau_{\mathcal X}$, by
\cite[Prop.~3.4.6]{Thu05}.  We use only the model-function case, in which
equality holds.
\end{remark}

\subsection{The slope formula}

The next theorem is the exact form in which the graph Laplacian will be used
in \cref{part:cld,part:bfj}.  We write $r_*$ for pushforward of divisors:
if $\operatorname{div}(f)=\sum_Pm_PP$ with $P$ closed points of $C$, then
\[
 r_*\operatorname{div}(f)
   =\sum_Pm_P\,[\kappa(P):K]\ \delta_{r(P)} ,
\]
a signed atomic measure on $\Gamma$ of total mass
$\deg\operatorname{div}(f)=0$.  The degree factor is invisible when $K$ is
algebraically closed and is forced in general by the residue-field degrees
occurring in intersection numbers.

The hypothesis in the next theorem is automatic when
$\operatorname{supp}\operatorname{div}(f)$ consists of $K$-rational
points: at a node with local equation $xy=\pi$, a point $P\in C(K)$ has
$v(x(P))\in v(K^\times)=\mathbf Z$ and $0\leq v(x(P))\leq1$, so
$v(x(P))\in\{0,1\}$ and $P$ does not reduce to the node.  In general it
holds after a finite extension of $K$; see
\cref{rem:ch08-slope-scope}.

\begin{theorem}
\label{thm:ch08-slope-formula}
Let $\cX$ be a regular semistable model of $C$ with skeleton $\Gamma$ and
retraction $r$, and let $f\in K(C)^\times$.  Assume that every point of
$\operatorname{supp}\operatorname{div}(f)$ reduces to a smooth point of
$\cX_s$; equivalently, that $r$ maps
$\operatorname{supp}\operatorname{div}(f)$ into the vertex set of $\Gamma$.
Write $F=v_\bullet(f)|_\Gamma$, that is, $F(p)=v_p(f)$ for $p\in\Gamma$.
Then $F$ is affine with integral slope on every edge of $\Gamma$, and
\begin{equation}
\label{eq:ch08-slope-formula}
 \Delta_\Gamma F=-\,r_*\operatorname{div}(f) ,
 \qquad\text{equivalently}\qquad
 \Delta_\Gamma\bigl(\log|f|\,\bigr|_\Gamma\bigr)=r_*\operatorname{div}(f) .
\end{equation}
\end{theorem}

\begin{proof}
\emph{Step 1: separating the reductions.}  We first show that both sides of
\eqref{eq:ch08-slope-formula} are unchanged if $\cX$ is replaced by the
blowup $\cX'$ of a smooth closed point $p$ of $\cX_s$.  Such a blowup
inserts an exceptional $E\simeq\mathbf P^1_{\widetilde K}$ with
$\ord_E(\pi)=\operatorname{mult}_p(\cX_s)=1$, meeting the strict transform of
the component through $p$ transversally at one point; so $\cX'$ is again
regular semistable, the hypothesis on reductions is preserved, and $\Gamma'$
is $\Gamma$ with one pendant edge $T$ attached at the vertex $p'$ carrying
$p$.  Moreover $\Gamma\subset\Gamma'$ and
$r_{\cX}=r_{\cX}\circ r_{\cX'}$ by \cref{sec:skeleton-refinement}, and
$r_\cX|_{\Gamma'}$ collapses $T$ to $p'$.

Write $F'=v_\bullet(f)|_{\Gamma'}$, so $F'|_\Gamma=F$.  The total mass of
$\Delta_{\Gamma'}F'$ on $T\setminus\{p'\}$ is
$-\sum_{\nu\in T_{p'}T}d_\nu F'(p')$, since the two contributions of every
edge interior to $T$ cancel and the edges issuing from $p'$ contribute only
at their far ends.  Hence
\[
 \begin{aligned}
 \bigl(r_\cX{}_*\Delta_{\Gamma'}F'\bigr)(\{p'\})
 &=\sum_{\nu\in T_{p'}\Gamma}d_\nu F'(p')
   +\sum_{\nu\in T_{p'}T}d_\nu F'(p')\\
 &\quad-\sum_{\nu\in T_{p'}T}d_\nu F'(p')
  =(\Delta_\Gamma F)(\{p'\}) ,
 \end{aligned}
\]
and at all other points of $\Gamma$ the two measures visibly agree; so
$r_\cX{}_*\Delta_{\Gamma'}F'=\Delta_\Gamma F$.  On the other side,
$r_\cX{}_*r_{\cX'}{}_*\operatorname{div}(f)=r_{\cX}{}_*\operatorname{div}(f)$.
Thus \eqref{eq:ch08-slope-formula} for $\cX'$ implies it for $\cX$.

Since $\operatorname{div}(f)$ is a finite set of closed points, all reducing
to smooth points of $\cX_s$ by hypothesis, finitely many such blowups make
their reductions pairwise distinct.  We assume this from now on.

\emph{Step 2: the intersection identity.}  Decompose the principal divisor of
$f$ on the model:
\[
 \operatorname{div}_{\cX}(f)=H+V,
 \qquad
 H=\overline{\operatorname{div}(f)},
 \qquad
 V=\sum_Yc_YY ,
\]
where $H$ is the horizontal part, the Zariski closure in $\cX$ of the divisor
of $f$ on $C$, and $V$ is vertical.  By definition
$c_Y=\ord_Y(f)=v_{v_Y}(f)=F(v_Y)$, the multiplicities $b_Y$ all being $1$.
Since $\operatorname{div}_\cX(f)$ is a principal Cartier divisor and every
component $Y$ of $\cX_s$ is a proper curve over $\widetilde K$,
\begin{equation}
\label{eq:ch08-principal-degree-zero}
 (H\cdot Y)+(V\cdot Y)=\deg\bigl(\cO_\cX(\operatorname{div}_\cX f)|_Y\bigr)=0 :
\end{equation}
the line bundle $\cO_\cX(\operatorname{div}_\cX f)$ is trivial, being
generated by $f^{-1}$, and the degree of a trivial line bundle on the proper
$\widetilde K$-curve $Y$ vanishes.  The identification of the intersection
numbers with residual degrees is
\cref{lem:vertical-intersection-degree}.

The same identity holds verbatim on any regular semistable model of $C_{K'}$
over the valuation ring of a finite extension $K'/K$, with $v'$ the extension
of $v$; only the normalization of the divisorial points changes, namely
$v_Z=v'(\varpi')\ord_Z$ for a uniformizer $\varpi'$ of $K'^\circ$ when
$\ord_Z(\varpi')=1$.

\emph{Step 3: $F$ is affine on every edge.}  Fix an edge $e$ of $\Gamma$,
attached to a node $\xi$ with formal equation $xy=\pi$, parametrized by
$s=v(x)\in[0,1]$ as in \cref{sec:ch01-running-nodal}.  Fix an integer
$q\geq1$, put $K_q=K(\pi^{1/q})$ and $\varpi=\pi^{1/q}$, so that the
extension $v_q$ of $v$ satisfies $v_q(\varpi)=1/q$.  Base changing the node
gives $xy=\varpi^q$, and by \cref{ex:snc-chain-models} the resolution of this
$A_{q-1}$-singularity is a regular semistable model over $K_q^\circ$ whose
special fibre over $e$ is a chain
\[
 E_0,E_1,\ldots,E_q ,
 \qquad
 (E_j\cdot E_{j\pm1})=1,\quad (E_j\cdot E_j)=-2\ \ (0<j<q),
\]
of components of multiplicity one, with $v_{E_j}(x)=j/q$.  The two ends
$E_0,E_q$ are the strict transforms of the two components of $\cX_s$ through
$\xi$, and the divisorial point of $E_j$ is $v_{E_j}=q^{-1}\ord_{E_j}$;
carrying this out for every node produces a regular semistable model $\cX_q$
of $C_{K_q}$ whose skeleton is again $\Gamma$, now subdivided into edges of
length $1/q$.

Because every point of $\operatorname{div}(f)$ reduces to a smooth point of
$\cX_s$, no point of $\operatorname{div}(f)$ lies in the open annulus over
$e^\circ$; hence the horizontal part $H_q$ over $K_q$ meets no $E_j$ with
$0<j<q$, so $(H_q\cdot E_j)=0$.  Applying Step~2 on $\cX_q$ with
$a_j=\ord_{E_j}(f)=qF(j/q)$ gives, for $0<j<q$,
\[
 0=(V_q\cdot E_j)=a_{j-1}-2a_j+a_{j+1}
  =q\Bigl(F\bigl(\tfrac{j-1}q\bigr)+F\bigl(\tfrac{j+1}q\bigr)
          -2F\bigl(\tfrac jq\bigr)\Bigr).
\]
Thus $j\mapsto F(j/q)$ is an arithmetic progression, that is,
$F(j/q)=F(0)+\tfrac jq\bigl(F(1)-F(0)\bigr)$ for $0\leq j\leq q$.  Since
$q$ was arbitrary, $F$ is affine on $\mathbf Q\cap[0,1]$, hence on $[0,1]$ by
continuity.  Its slope $F(1)-F(0)=\ord_{Y'}(f)-\ord_Y(f)$ is an integer.

\emph{Step 4: the two terms.}  By Step~1 the points of
$\operatorname{supp}\operatorname{div}(f)$ have distinct smooth reductions.
The closure $\overline{\{P\}}$ is finite and flat over $R$ of degree
$[\kappa(P):K]$, so its special fibre has length $[\kappa(P):K]$ over
$\widetilde K$ and is concentrated at the reduction of $P$; that reduction
lies on $Y$ exactly when $r(P)=v_Y$.  Hence
\[
 (H\cdot Y)=\sum_{P:\,r(P)=v_Y}m_P\,[\kappa(P):K]
          =\bigl(r_*\operatorname{div}(f)\bigr)(\{v_Y\}) .
\]
For the vertical term, use that $\cX_s=\operatorname{div}(\pi)$ is principal,
so that $(\cX_s\cdot Y)=0$; since $\cX_s=\sum_{Y'}Y'$ with all
multiplicities one, this reads
$(Y\cdot Y)=-\sum_{Y'\neq Y}(Y\cdot Y')$.  The number $(Y\cdot Y')$ is the
number of nodes lying on $Y\cap Y'$, that is, the number of edges of $\Gamma$
joining $v_Y$ to $v_{Y'}$.  Therefore
\[
 (V\cdot Y)=\sum_{Y'}c_{Y'}(Y'\cdot Y)
          =\sum_{Y'\neq Y}(c_{Y'}-c_Y)(Y\cdot Y')
          =\sum_{e\ni v_Y}\bigl(F(\text{far end of }e)-F(v_Y)\bigr).
\]
Each edge has length one and $F$ is affine on it by Step~3, so the summand is
precisely the outgoing slope $d_eF(v_Y)$, and
$(V\cdot Y)=(\Delta_\Gamma F)(\{v_Y\})$.  This computation is a special case
of the degree bookkeeping recorded in \cref{cor:special-fibre-relation}.

\emph{Step 5: conclusion.}  Substituting Step~4 into
\eqref{eq:ch08-principal-degree-zero} gives, for every component $Y$,
\[
 (\Delta_\Gamma F)(\{v_Y\})
 =-(r_*\operatorname{div}f)(\{v_Y\}).
\]
By Step~3 the measure $\Delta_\Gamma F$ is supported on the
vertices, and by hypothesis so is $r_*\operatorname{div}(f)$.  This proves
the first identity in \eqref{eq:ch08-slope-formula}, and the second follows
from $\log|f|=-v_\bullet(f)$ and linearity of $\Delta_\Gamma$.
\end{proof}

\begin{remark}[Scope of the hypothesis]
\label{rem:ch08-slope-scope}
Suppose a point $P$ of $\operatorname{supp}\operatorname{div}(f)$ reduces to
a node $\xi$ with equation $xy=\pi$.  Then $c=v(x(P))$ lies in
$(0,1)\cap\mathbf Q$, because $\kappa(P)$ is finite over $K$; write
$c=p/q$.  Base changing to $K_q=K(\pi^{1/q})$ and resolving as in Step~3
produces a regular semistable model of $C_{K_q}$ on which the points of
$C_{K_q}$ above $P$ reduce into the smooth locus of the component $E_{p}$,
because they satisfy $v_q(x)=p/q$ exactly.  Doing this for the finitely many
offending points at once, the hypothesis of
\cref{thm:ch08-slope-formula} is satisfied over a finite extension of $K$.
Neither side of \eqref{eq:ch08-slope-formula} changes: the extension $v_q$ of
$v$ leaves $F$ unchanged, and $r_*\operatorname{div}(f)$ is unchanged because
the degree $[\kappa(P):K]$ of a closed point is distributed among the points
above it.  Consequently \eqref{eq:ch08-slope-formula} holds without the
displayed hypothesis, atoms of $r_*\operatorname{div}(f)$ being then allowed
in the interiors of the edges of $\Gamma$; we have kept the hypothesis in the
statement because every application in this book satisfies it, and because
with it the proof is entirely intersection-theoretic.
\end{remark}

\begin{corollary}
\label{cor:ch08-harmonic-away}
With the notation of \cref{thm:ch08-slope-formula}, the function
$\log|f|$ restricted to $\Gamma$ is harmonic outside
$r(\operatorname{supp}\operatorname{div}(f))$.  If $f$ has neither zeros nor
poles on $C$, then $f$ is constant and $F$ is constant.
\end{corollary}

\begin{proof}
The first assertion is \eqref{eq:ch08-slope-formula}.  For the second, a
regular invertible function on a projective curve is constant, so
$\operatorname{div}(f)=0$; then $\Delta_\Gamma F=0$ and
\cref{lem:ch08-laplacian-well-defined}\,(3) applies.
\end{proof}

\begin{remark}[Scope]
\Cref{thm:ch08-slope-formula} is stated for a regular semistable model, which
is all that is used later, and the proof is complete under that hypothesis.
The same formula, with multiplicities $b_Y$ inserted in the normalization of
the divisorial points and in the edge lengths, holds on the skeleton of an
arbitrary SNC model; the harmonic-analysis version valid on all of
$C^{\An}$, including the type-$\mathrm{IV}$ points and the branches off the
skeleton, is Thuillier's potential theory on Berkovich curves \cite{Thu05}
and is compared with the Chambert--Loir--Ducros calculus in
\cref{chap:cld-examples}, see \cref{thm:wanner-curve-comparison} and
\cite{Wanner19,GR25PL}.  We do not use that stronger statement in
\cref{part:foundations}.

Concretely, the two ingredients \cref{thm:ch08-slope-formula} extracts from
intersection theory are proved sheaf-theoretically in \cite{Thu05}: that
$\log|f|$ is harmonic wherever $f$ is an invertible section of $\cO_C$
\cite[Prop.~2.3.20]{Thu05}, and that it is subharmonic wherever $f$ is a
section invertible on a dense open \cite[Prop.~3.1.6]{Thu05}.  The atoms are
then produced by the Poincar\'e--Lelong formula \cite[\S3.3.3]{Thu05}.
\end{remark}

\begin{remark}[Harmonic functions and $\log|\cO^\times|$ are not the same
sheaf]
\label{rem:ch08-harmonic-vs-log}
It is tempting to \emph{define} harmonicity on a curve by ``locally of the
form $\log|f|$ with $f$ invertible''.  Thuillier shows that this fails, and
identifies the obstruction exactly.  Write $\mathcal F_C$ for the sheaf of
$\mathbf R$-vector spaces on $|C^{\An}|$ generated by the germs $\log|f|$,
$f\in\Gamma(U,\cO_{C^{\An}}^{\times})$, and $\mathcal H_C$ for his sheaf of
harmonic germs.  Then $\mathcal F_C\subseteq\mathcal H_C$, the quotient is a
skyscraper sheaf supported on the type-$\mathrm{II}$ points, and at such a
point $x$ with associated residue curve $C_x$ over $\widetilde K$ one has
\[
 (\mathcal H_C/\mathcal F_C)_x\simeq\Pic^0(C_x)\otimes_{\mathbf Z}\mathbf R ,
\]
by \cite[Lem.~2.3.22]{Thu05}: an $\mathbf R$-divisor of degree zero on a
component of the special fibre is the divisor of a harmonic function, but
need not be principal.  The quotient vanishes---so that the naive definition
is correct after all---if $\widetilde K$ is algebraic over a finite field, or
if $C^{\An}$ is locally isomorphic to $\mathbf P^{1,\An}$, that is if $C$ is
a Mumford curve; this is \cite[Thm.~2.3.21]{Thu05}.  The same two hypotheses
reappear, for the same reason, in Wanner's earlier ordinary-smooth
regularization theorem.  The positive-curvature construction of
\Cref{thm:p063-current-psh-ordinary-regularization}, in its scope of an
algebraically closed complete nontrivially valued base field, removes those
curve and residue-field restrictions for psh, but not harmonic,
approximants.
\end{remark}

\subsection{Model measures on a curve as graph data}
\label{sec:ch08-curve-cl}

The slope formula makes the Chambert--Loir measure of a model metric on a
curve into a purely graph-theoretic object.

\begin{proposition}
\label{prop:ch08-curve-cl}
Let $\cX$ be a regular semistable model of $C$, let $L$ be a line bundle on
$C$ of degree $d$, and let $\cL$ be a model of $L$ on $\cX$, with associated
model metric $\|\cdot\|_{\cL}$ and skeleton $\Gamma=\Sk(\cX)$.  Then
\begin{equation}
\label{eq:ch08-curve-cl-measure}
 c_1(L,\|\cdot\|_{\cL})=\sum_Y\deg(\cL|_Y)\,\delta_{v_Y},
 \qquad
 \sum_Y\deg(\cL|_Y)=d .
\end{equation}
If $\cL'=\cL\otimes\cO_\cX(D)$ for a vertical divisor $D$ with model function
$\varphi_D$, then
\begin{equation}
\label{eq:ch08-curve-cl-perturbed}
 c_1(L,\|\cdot\|_{\cL'})
   =c_1(L,\|\cdot\|_{\cL})+\Delta_\Gamma\bigl(\varphi_D|_\Gamma\bigr),
\end{equation}
and the metric $\|\cdot\|_{\cL'}$ is semipositive if and only if the
right-hand side is a positive measure.
\end{proposition}

\begin{proof}
The first formula is the definition of the model measure in
\cref{sec:metrics-cl-measure} with all $b_Y=1$, and the mass statement is
\cref{thm:cl-model-properties}.  For the second, note that $\varphi_D$ is
affine on every edge of $\Gamma$ after subdivision
(\cref{prop:skeleton-model-pa}), so that $\Delta_\Gamma\varphi_D$ makes
sense; and by the computation of Step~4 in the proof of
\cref{thm:ch08-slope-formula}, applied to the vertical divisor
$D=\sum_Ya_YY$ instead of $V$,
\[
 \deg\bigl(\cO_{\cX}(D)|_Y\bigr)=(D\cdot Y)
 =\sum_{e\ni v_Y}\bigl(\varphi_D(\text{far end of }e)-\varphi_D(v_Y)\bigr)
 =(\Delta_\Gamma\varphi_D)(\{v_Y\}) ,
\]
using $\varphi_D(v_Y)=a_Y$.  If $D$ is not determined on $\cX$, pass to a
dominating regular semistable model, on which the computation applies and on
which the measure of $\cL$ is unchanged by
\cref{thm:cl-model-properties}; the additional atoms of
$\Delta_\Gamma\varphi_D$ then sit at the new vertices, in the interior of the
old edges.  Additivity of degrees gives
\eqref{eq:ch08-curve-cl-perturbed}.  Finally, a model metric on a curve is
semipositive precisely when all the residual degrees are nonnegative
(\cref{sec:metrics-semipositive}), which is positivity of the measure.
\end{proof}

Formula \eqref{eq:ch08-curve-cl-perturbed} is exactly
\eqref{eq:graph-curvature} with $\mu_0=c_1(L,\|\cdot\|_\cL)$: the model
picture and the graph picture agree, and the mass $d$ is conserved because
$\Delta_\Gamma$ kills total mass by \eqref{eq:lap-mass-zero}.  All the
curve computations of this chapter are instances of this one identity.

\section{The analytic projective line}
\label{sec:berkovich-p1}

For the rest of this section $K$ is complete and nontrivially valued, with
value group $v(K^\times)\subseteq\mathbf R$ and residue field $\widetilde K$.
Two additional hypotheses are used in disjoint places and are never used
together, since they are incompatible.  In
\cref{sec:ch08-p1-tree} we assume $K$ algebraically closed, as
\cref{thm:berk-a1-classification} requires; the reader may keep
$K=\mathbf C_p$ in mind.  In
\crefrange{sec:ch08-p1-newton}{sec:ch08-p1-models}, where intersection
numbers on models are computed, we assume instead that $K$ is discretely
valued with $v(\pi)=1$ and $R=K^\circ$; there the reader may keep
$K=\mathbf Q_p$ or $K=k((\pi))$ in mind.

\subsection{The tree and its metric}
\label{sec:ch08-p1-tree}
\index{Berkovich projective line!tree}\index{path metric}

\Cref{thm:berk-a1-classification} classifies the points of
$\mathbf A^{1,\An}_K$ into four types and
\cref{prop:berk-a1-tree} organizes them into a tree.  We add the metric and
the branching data, which is what potential theory needs.

\begin{definition}[Path metric]
\label{def:ch08-path-metric}
For a point $x$ of type $\mathrm{II}$ or $\mathrm{III}$, corresponding to
the disc $D(a,r)$, put
\[
                       t(x)=-\log r\in\mathbf R .
\]
If $x\leq y$ are two such points, the segment $[x,y]$ is the totally ordered
set $\{z:x\leq z\leq y\}$, which by \cref{prop:berk-a1-tree}\,(1) is the
family of Gauss points $\zeta_{a,\rho}$ with $r\leq\rho\leq s$; we give it the
length $|t(x)-t(y)|$.  For arbitrary $x,y$ of type
$\mathrm{II}$ or $\mathrm{III}$ the path from $x$ to $y$ is
$[x,x\vee y]\cup[x\vee y,y]$, and
\begin{equation}
\label{eq:ch08-path-metric}
 \rho(x,y)=\bigl(t(x)-t(x\vee y)\bigr)+\bigl(t(y)-t(x\vee y)\bigr) .
\end{equation}
\end{definition}

Thus $t$ decreases as discs grow, in accordance with $v=-\log|\cdot|$: the
Gauss point of the unit disc has $t=0$, the point $\zeta_{0,|\pi|^m}$ has
$t=m$, and $t\to-\infty$ approaches the point $\infty\in\mathbf P^1$.  This
is the coordinate already used in \eqref{eq:gauss-disc} below and in
\cref{sec:ch01-running-toric}.

\begin{example}[A distance computed]
\label{ex:ch08-path-distance}
Suppose $p\in K^\times$ has $v(p)=1$ and take
$x=\zeta_{0,|p|^3}$ and $y=\zeta_{1,|p|^5}$.  Then $|0-1|=1$, so
\cref{prop:berk-a1-tree}\,(3) gives
$x\vee y=\zeta_{0,\max\{|p|^3,|p|^5,1\}}=\zeta_{0,1}$, the Gauss point, with
$t=0$.  Hence $\rho(x,y)=3+5=8$.  The two points lie on different branches
issuing from $\zeta_{0,1}$, and every path between them passes through it.
\end{example}

\begin{proposition}
\label{prop:ch08-branching}
Let $x=\zeta_{a,r}$ with $r\in|K^\times|$, a point of type $\mathrm{II}$.
The set of tangent directions at $x$ in $\mathbf P^{1,\An}_K$ is in
canonical bijection with $\mathbf P^1(\widetilde K)$.  If $r\notin|K^\times|$,
that is, if $x$ is of type $\mathrm{III}$, there are exactly two tangent
directions.
\end{proposition}

\begin{proof}
After the change of variable $T\mapsto(T-a)/c$ with $|c|=r$, which is an
automorphism of $\mathbf P^{1,\An}_K$ preserving types and the tree
structure, we may assume $a=0$ and $r=1$, so $x=\zeta_{0,1}$ is the Gauss
point of the unit disc.  A tangent direction at $x$ is an equivalence class
of points $y\neq x$, two points being equivalent when the segments $[x,y]$
and $[x,y']$ share an initial segment; equivalently, it is a connected
component of $\mathbf P^{1,\An}_K\setminus\{x\}$.

The components are described by the reduction map
$\mathrm{red}:\mathbf P^{1,\An}_K\to\mathbf P^1_{\widetilde K}$ attached to
the model $\mathbf P^1_{K^\circ}$ in \cref{sec:berk-reduction}, whose fibre
over the generic point of the special fibre is $\{x\}$ by
\cref{prop:berk-divisorial-fibre}.  If
$y\neq x$ and $y$ lies in the closed unit disc, then $y\leq\zeta_{0,1}$ and
$y$ determines a disc $D(b,\rho)$ with $|b|\leq1$ and $\rho<1$; two such
points lie in the same component of the complement of $x$ exactly when their
discs have a common point, that is, when the residues
$\widetilde b\in\widetilde K$ coincide.  This gives one direction for each
element of $\widetilde K$, and the remaining component, containing all points
with $|T|>1$ and the point $\infty$, gives the direction of
$\infty\in\mathbf P^1(\widetilde K)$.  Since $|K^\times|$ is the value group
and $r=1$ lies in it, all residues $\widetilde b$ are realized.

If $r\notin|K^\times|$, the same argument produces only two components: the
points strictly below $x$ and the points not below $x$.  Indeed, if
$y<x$ corresponds to $D(b,\rho)$ with $\rho<r$, and $y'<x$ corresponds to
$D(b',\rho')$ with $\rho'<r$, then $|b-b'|\leq r$ and $|b-b'|\neq r$ because
$r\notin|K^\times|$, so $|b-b'|<r$ and the two discs are contained in a
common disc of radius $<r$; hence $y$ and $y'$ lie in the same component.
\end{proof}

\begin{remark}[Where the branches go]
\Cref{prop:ch08-branching} is the analytic form of the statement that
blowing up a closed point of the special fibre of $\mathbf P^1_R$ creates a
new component: the closed points of the special fibre
$\mathbf P^1_{\widetilde K}$ are exactly $\mathbf P^1(\widetilde K)$, and each
of them is the direction of a branch at the Gauss point.  We compute two such
blowups in \cref{ex:ch08-p1-tripod}.
\end{remark}

\subsection{Gauss points, tropical coordinates, and Newton polygons}
\label{sec:ch08-p1-newton}

Let $K$ be complete and nontrivially valued.  For $r>0$, the Gauss point
$\zeta_{0,r}$ of the closed disc centered at $0$ is defined by
\begin{equation}
\label{eq:gauss-disc}
 \left|\sum_{i=0}^m a_iT^i\right|_{\zeta_{0,r}}
 =\max_i|a_i|r^i.
\end{equation}
Put $t=-\log r$.  Then
\[
                    v_{\zeta_{0,r}}(T)=t,
\qquad
 v_{\zeta_{0,r}}\left(\sum_i a_iT^i\right)
 =\min_i\{v(a_i)+it\}.
\]
The points $\zeta_{0,r}$ form the skeleton of the analytic torus
$\mathbf G_{m,K}^{\mathrm{an}}$, canonically identified with $\mathbf R$ via
$t$.  Its compactification in $(\mathbf P^1_K)^{\mathrm{an}}$ has endpoints
$0$ and $\infty$.  The entire Berkovich projective line is a real tree
obtained by attaching further branches to this apartment.

\begin{example}[A polynomial along the skeleton]
For $f(T)=a_0+a_1T+\cdots+a_mT^m$,
\[
 -\log|f(\zeta_{0,\rme^{-t}})|
 =\min_i\{v(a_i)+it\}.
\]
This is a concave piecewise affine function with integral slopes.  A slope
changes precisely when two terms have the same Gauss size.  Thus the Newton
polygon of $f$ is visible as the restriction of $v(f)$ to the skeleton.
\end{example}

\begin{example}[A cubic, its Newton polygon, and the slope formula]
\label{ex:ch08-newton-cubic}
\index{Newton polygon}
Let $v(p)=1$ and consider
\[
                 f(T)=T^3+p^2T^2+p^5T+p^9 .
\]
Along the skeleton, with $t=v_t(T)$,
\[
 F(t):=v_t(f)=\min\{3t,\ 2+2t,\ 5+t,\ 9\} .
\]
The four affine functions have slopes $3,2,1,0$ and are compared pairwise:
$3t=2+2t$ at $t=2$; $2+2t=5+t$ at $t=3$; $5+t=9$ at $t=4$.  A direct check
at the three candidate breakpoints confirms that the minimum is attained by
consecutive terms only:
\[
 F(t)=
 \begin{cases}
  3t, & t\leq2,\\
  2+2t, & 2\leq t\leq3,\\
  5+t, & 3\leq t\leq4,\\
  9, & t\geq4,
 \end{cases}
 \qquad
 F(2)=6,\quad F(3)=8,\quad F(4)=9 .
\]
Thus $F$ is concave with slopes $3,2,1,0$ and slope jumps of size $-1$ at
$t=2,3,4$.  Each of the three segments of the Newton polygon has horizontal
length one, so $f$ has exactly one root of valuation $2$, one of valuation
$3$, and one of valuation $4$, and each of them is a $K$-rational point of
$\mathbf P^1$; call them $P_2,P_3,P_4$.  (If one prefers not to invoke the
Newton-polygon factorization, pass to an algebraic closure: neither side of
the identity below changes.)

We now verify \cref{thm:ch08-slope-formula} by hand on the compact skeleton
$\Gamma=[\zeta_{0,1},\zeta_{0,|p|^4}]$, that is, on $0\leq t\leq4$; this
$\Gamma$ is the skeleton of the model of $\mathbf P^1$ obtained from
$\mathbf P^1_R$ by four successive blowups of the point $T=0$ of the special
fibre, whose chain of exceptional curves has divisorial points at
$t=1,2,3,4$.  Compute the outgoing-slope sums of $F$ on $\Gamma$:
\begin{align*}
 t=0:&\quad d_\nu F=+3
   &&\Longrightarrow\quad \text{mass }+3;\\
 t=1:&\quad d_\nu F=+3\ \text{ and }\ -3
   &&\Longrightarrow\quad \text{mass }\ \ 0;\\
 t=2:&\quad d_\nu F=+2\ \text{ and }\ -3
   &&\Longrightarrow\quad \text{mass }-1;\\
 t=3:&\quad d_\nu F=+1\ \text{ and }\ -2
   &&\Longrightarrow\quad \text{mass }-1;\\
 t=4:&\quad d_\nu F=-1
   &&\Longrightarrow\quad \text{mass }-1 .
\end{align*}
Hence
\[
 \Delta_\Gamma F=3\delta_{t=0}-\delta_{t=2}-\delta_{t=3}-\delta_{t=4},
\]
of total mass $3-1-1-1=0$, as \eqref{eq:lap-mass-zero} demands.  On the other
side, $\operatorname{div}(f)=P_2+P_3+P_4-3\,(\infty)$ on $\mathbf P^1$, and
the retraction onto $\Gamma$ sends $P_j$ to the point $t=j$ for $j=2,3,4$
and sends $\infty$, which lies outside the unit disc, to the endpoint
$t=0$.  Therefore
\[
 r_*\operatorname{div}(f)
 =\delta_{t=2}+\delta_{t=3}+\delta_{t=4}-3\delta_{t=0}
 =-\Delta_\Gamma F ,
\]
which is \eqref{eq:ch08-slope-formula}.  Note where each ingredient of the
proof appears: concavity of $F$ is the ultrametric inequality, the integrality
of the slopes is the integrality of $\ord$, and the mass $+3$ at the endpoint
$t=0$ is the pole of order three at infinity, pushed onto the skeleton.
\end{example}

For a nonzero rational function $f$ on $\mathbf P^1$, the slope formula is
the Poincar\'e--Lelong identity
\begin{equation}
\label{eq:p1-poincare-lelong}
                 \dprimeddprime[\log|f|]=\delta_{\operatorname{div}(f)}.
\end{equation}
In particular,
\[
                 \dprimeddprime[\log|T|]=\delta_0-\delta_\infty.
\]
Since $v(T)=-\log|T|$, its Laplacian has the opposite sign.  This check will
prevent sign changes when we pass between valuations and metric weights.

\subsection{The canonical metric on the hyperplane bundle}

Let $[Z_0:Z_1]$ be homogeneous coordinates and $T=Z_1/Z_0$.  The standard
model metric on $\mathcal O(1)$ satisfies
\[
 \|Z_0(x)\|=\frac{|Z_0(x)|}{\max\{|Z_0(x)|,|Z_1(x)|\}}.
\]
On the chart $Z_0\neq0$, its weight is
\begin{equation}
\label{eq:p1-standard-weight}
      -\log\|Z_0\|(x)=\log\max\{1,|T(x)|\}.
\end{equation}
Along the toric skeleton this becomes
\[
                         F(t)=\max\{0,-t\}.
\]
It has one unit slope jump, at the Gauss point $\zeta_{0,1}$.  Therefore
\begin{equation}
\label{eq:p1-standard-measure}
 c_1\bigl(\mathcal O(1),\|\cdot\|_{\mathrm{mod}}\bigr)
 =\delta_{\zeta_{0,1}}.
\end{equation}
The total mass is $1=\deg\mathcal O(1)$.  For $\mathcal O(d)$ with the tensor
power metric, the measure is $d\delta_{\zeta_{0,1}}$.

\begin{remark}
Formula \eqref{eq:p1-standard-measure} may look singular from the complex
viewpoint: all curvature is concentrated at one nonclassical point.  It is
nevertheless the exact non-Archimedean analogue of a semipositive model
metric.  Smoothness will later refer to smoothness on tropical charts, not
to diffuseness of the curvature measure.
\end{remark}

\subsection{Models of the projective line and their Laplacians}
\label{sec:ch08-p1-models}
\index{model!of the projective line}

Assume now that $K$ is discretely valued, $R=K^\circ$, $v(\pi)=1$.  The
simplest model of $X=\mathbf P^1_K$ is $\cX_0=\mathbf P^1_R$; its special
fibre is the single reduced component $\mathbf P^1_{\widetilde K}$, its
skeleton is the single point $\zeta_{0,1}$, and its retraction sends
everything to the Gauss point.  All other models used here are obtained from
$\cX_0$ by blowing up closed points of special fibres: the possible centres
on $\cX_{0,s}$ are the closed points of $\mathbf P^1_{\widetilde K}$, which
are indexed by $\mathbf P^1(\widetilde K)$ exactly when $\widetilde K$ is
algebraically closed (\cref{prop:ch08-branching}); for
$\widetilde K=\mathbf F_p$ there are infinitely many, one for each monic
irreducible polynomial.  In every case each blowup extends the skeleton by
one edge of length one, since the exceptional divisor has $b_E=1$.

We compute three models completely.  In each case the curvature measure of a
model metric on $\cO(1)$ is obtained twice, from intersection numbers by
\eqref{eq:ch08-curve-cl-measure} and from outgoing slopes by
\eqref{eq:ch08-curve-cl-perturbed}, and the two answers are compared.

\begin{example}[One blowup: the interval]
\label{ex:ch08-p1-interval}
Let $\cX_1\to\cX_0$ be the blowup of the point $\{T=\pi=0\}$, as in
\cref{ex:ch01-blowup-model-function}.  Its special fibre is
$\cX_{1,s}=\widetilde C+E$, both reduced, meeting transversally at one point,
and
\[
 (E\cdot\widetilde C)=1,\qquad (E\cdot E)=-1,\qquad
 (\widetilde C\cdot\widetilde C)=-1 ,
\]
the last two because $\cX_{1,s}=\operatorname{div}(\pi)$ has degree zero on
each component.  The skeleton is the segment
$\Gamma_1=[\zeta_{0,1},\zeta_{0,|\pi|}]$, of length one, with $t=0$ at
$x_{\widetilde C}$ and $t=1$ at $x_E$, and
$\varphi_E|_{\Gamma_1}(t)=t$ by \eqref{eq:ch01-blowup-model-function}.

Let $\cL_a=\rho^*\cO_{\cX_0}(1)\otimes\cO_{\cX_1}(aE)$ for $a\in\mathbf Z$,
where $\rho:\cX_1\to\cX_0$ is the blowup.  Its residual degrees are
\[
 \deg(\cL_a|_{\widetilde C})=1+a(E\cdot\widetilde C)=1+a,
 \qquad
 \deg(\cL_a|_{E})=0+a(E\cdot E)=-a ,
\]
using $\deg(\rho^*\cO(1)|_E)=0$ because $\rho$ contracts $E$.  Hence
\begin{equation}
\label{eq:ch08-p1-interval-measure}
 c_1\bigl(\cO(1),\|\cdot\|_{\cL_a}\bigr)
   =(1+a)\,\delta_{x_{\widetilde C}}-a\,\delta_{x_E},
\end{equation}
of total mass $(1+a)+(-a)=1=\deg\cO(1)$.  The metric is semipositive exactly
when $-1\leq a\leq0$.  The graph computation agrees: the model function of
$aE$ restricted to $\Gamma_1$ is $at$, whose outgoing-slope sums are $a$ at
$t=0$ and $-a$ at $t=1$, so
\[
 \delta_{x_{\widetilde C}}+\Delta_{\Gamma_1}(at)
 =(1+a)\delta_{x_{\widetilde C}}-a\delta_{x_E},
\]
matching \eqref{eq:ch08-p1-interval-measure}.  For $a=-1$ the whole degree
has migrated to the new component: the measure is $\delta_{x_E}$, and the
corresponding metric is the standard model metric transported by the
automorphism $T\mapsto\pi^{-1}T$.
\end{example}

\begin{example}[Two blowups: a two-legged tree]
\label{ex:ch08-p1-tripod}
\index{tree!two-legged model skeleton}
Blow up $\cX_0=\mathbf P^1_R$ at the two closed points
$\{T=\pi=0\}$ and $\{T-1=\pi=0\}$ of the special fibre; they are distinct
because their residues $0$ and $1$ are distinct in
$\mathbf P^1(\widetilde K)$.  Call the result $\cX_2$, with exceptional
curves $E_0$ and $E_1$
and strict transform $\widetilde C$.  Then
\[
 (E_i\cdot\widetilde C)=1,\quad (E_i\cdot E_i)=-1,\quad
 (E_0\cdot E_1)=0,\quad (\widetilde C\cdot\widetilde C)=-2 .
\]
The skeleton is a path of length two: its central vertex
$x_{\widetilde C}=\zeta_{0,1}$ has two legs of length one ending at
$x_{E_0}=\zeta_{0,|\pi|}$ and $x_{E_1}=\zeta_{1,|\pi|}$.  Note that the two
legs really are distinct branches at $\zeta_{0,1}$, by
\cref{prop:ch08-branching}, since $\widetilde0\neq\widetilde1$; and
$\rho(x_{E_0},x_{E_1})=2$ by \eqref{eq:ch08-path-metric}, since
$x_{E_0}\vee x_{E_1}=\zeta_{0,1}$.

Take $D=a_0E_0+a_1E_1$ and
$\cL=\rho^*\cO_{\cX_0}(1)\otimes\cO_{\cX_2}(D)$.  The residual degrees are
\[
 \deg(\cL|_{\widetilde C})=1+a_0+a_1,
 \qquad
 \deg(\cL|_{E_0})=-a_0,
 \qquad
 \deg(\cL|_{E_1})=-a_1 ,
\]
so that
\begin{equation}
\label{eq:ch08-p1-tripod-measure}
 c_1\bigl(\cO(1),\|\cdot\|_{\cL}\bigr)
 =(1+a_0+a_1)\delta_{x_{\widetilde C}}
  -a_0\delta_{x_{E_0}}-a_1\delta_{x_{E_1}} ,
\end{equation}
of total mass $1$.  Semipositivity holds exactly when $a_0\leq0$,
$a_1\leq0$, and $a_0+a_1\geq-1$: a triangle in the $(a_0,a_1)$-plane, which
is the nef cone of this model in the notation of
\cref{sec:metrics-semipositive}.

The graph computation again agrees.  On this path, $\varphi_{E_i}$ is the
function equal to the distance from the central vertex along the $i$-th leg
and to $0$ elsewhere, so $\varphi_D$ takes the value $0$ at the centre and
$a_i$ at the tip of leg $i$, and is affine on each leg.  The outgoing-slope
sums are $a_0+a_1$ at the centre and $-a_i$ at the tip of leg $i$, so
\[
 \delta_{x_{\widetilde C}}+\Delta_\Gamma\varphi_D
 =(1+a_0+a_1)\delta_{x_{\widetilde C}}
   -a_0\delta_{x_{E_0}}-a_1\delta_{x_{E_1}} ,
\]
as in \eqref{eq:ch08-p1-tripod-measure}.  A convenient symmetric semipositive
choice is $a_0=a_1=-\tfrac12$, allowed after passing to
$\mathbf Q$-model metrics: the measure becomes
\[
 0\cdot\delta_{x_{\widetilde C}}
 +\tfrac12\delta_{x_{E_0}}+\tfrac12\delta_{x_{E_1}} ,
\]
of mass $1$, with no mass at the Gauss point.  The corresponding metric on
$\cO(1)$ has weight, in the frame $Z_0$,
\[
 \log\max\{1,|T|\}
   -\tfrac12\min\bigl\{\max\{v(T),0\},1\bigr\}
   -\tfrac12\min\bigl\{\max\{v(T-1),0\},1\bigr\},
\]
by \eqref{eq:p1-standard-weight} and
\eqref{eq:ch01-blowup-model-function}.  It is the exact non-Archimedean
analogue of a metric whose curvature is half the equilibrium measure of one
small disc plus half that of another.
\end{example}

\begin{example}[A Fubini--Study metric and a moving atom]
\label{ex:ch08-p1-fs}
Fix $\lambda\in v(K^\times)$ and consider the Fubini--Study metric on
$\cO(1)$ of \cref{sec:metrics-fs} attached to the two sections $Z_0,Z_1$
with weights $0$ and $-\lambda$, that is, the metric with weight
\[
 \phi_\lambda=\max\{0,\ \log|T|-\lambda\}
\]
in the frame $Z_0$.  On the skeleton, with $t=v(T)$, this is
$\max\{0,-t-\lambda\}$, a convex piecewise affine function with a single
corner at $t=-\lambda$ and slopes $-1$ and $0$.  Its outgoing-slope sum at
the corner is $(+1)+0=1$ and vanishes elsewhere, so
\[
 c_1\bigl(\cO(1),\|\cdot\|_{\phi_\lambda}\bigr)=\delta_{\zeta_{0,\rme^{\lambda}}} ,
\]
a single Dirac mass of mass $1=\deg\cO(1)$ at the Gauss point of the disc of
radius $\rme^{\lambda}$.  As $\lambda$ varies over $v(K^\times)$, the atom
sweeps out the integral-model points of the apartment; allowing
$\lambda\in\mathbf Qv(K^\times)$ and \(\mathbf Q\)-model metrics sweeps out
all rational divisorial points.  For $\lambda=0$ we
recover \eqref{eq:p1-standard-measure}.
\end{example}

\begin{example}[A degree-two metric with two atoms]
\label{ex:ch08-p1-degree-two}
Let $L=\cO(2)$, of degree $2$, and take the metric with weight
\[
 \psi=\max\{0,-t\}+\max\{0,-t-1\}
\]
on the skeleton, in the frame $Z_0^{\otimes2}$; it is the tensor product of
the standard model metric with the metric $\phi_1$ of
\cref{ex:ch08-p1-fs}.  Its Laplacian is
\[
 \Delta\psi=\delta_{t=0}+\delta_{t=-1},
\]
so
\[
 c_1\bigl(\cO(2),\|\cdot\|_\psi\bigr)
   =\delta_{\zeta_{0,1}}+\delta_{\zeta_{0,|\pi|^{-1}}} ,
\]
of total mass $2=\deg\cO(2)$, as required by
\cref{thm:cl-model-properties}.  Contrast this with the $d$-th tensor power
of a single model metric, whose measure is $d\delta_{\zeta_{0,1}}$: the same
mass may be distributed over one atom or several, and the possible
distributions are exactly the positive measures of mass $d$ obtained from
\cref{prop:graph-poisson}.
\end{example}

\begin{example}[Mass conservation under refinement]
\label{ex:ch08-p1-refinement}
Pull back the standard model metric on $\cO(1)$ from $\cX_0$ to the chain
model $\cX_4$ with special fibre $\widetilde C+E_1+E_2+E_3+E_4$ used in
\cref{ex:ch08-newton-cubic}.  Since $\rho$ contracts every $E_j$, the
residual degrees are $1,0,0,0,0$, and the measure is still
$\delta_{\zeta_{0,1}}$.  This is \cref{thm:cl-model-properties} in the
smallest instance: passing to a higher model neither creates nor destroys
mass, it only creates points at which mass \emph{could} sit.  The reason the
point $t=1$ contributed nothing in \cref{ex:ch08-newton-cubic} is different,
and local to that example: there $F$ is affine across $t=1$, the two
outgoing slopes being $+3$ and $-3$, because $f$ has no root of valuation
one.
\end{example}

\section{The trivially valued cone}
\label{sec:trivial-valuation-cone}

Over a trivially valued ground field the picture changes in one decisive
respect: there is no element $\pi$ with $v(\pi)=1$, hence no distinguished
slice of the space of semivaluations, and the scaling action
$v\mapsto tv$ becomes part of the structure rather than a normalization to be
fixed.  This section computes the resulting cone in the two smallest cases,
a smooth curve and the affine plane, and computes the homogeneous potentials
that \cref{part:bj} takes as its basic objects.  The general definitions are
in \cref{chap:trivial-analytification}
(\cref{def:trivial-analytification,def:normalized-link}) and
\cref{chap:bj-homogeneity-bdivisors}
(\cref{def:homogeneous-function}); nothing below presupposes them.

\subsection{The cone over a smooth curve}
\label{sec:ch08-trivial-curve}

Let $k$ be algebraically closed with the trivial absolute value.  On
$\mathbf A^1_k$ and for $a\in k$, define
\[
 v_{a,t}(f)=t\operatorname{ord}_a(f),\qquad 0\leq t<+\infty.
\]
At $t=0$ this is the trivial valuation.  There is also the endpoint
semivaluation
\[
 v_{a,\infty}(f)=
 \begin{cases}
  +\infty,&f(a)=0,\\
  0,&f(a)\neq0,
 \end{cases}
\]
whose support is the maximal ideal of $a$.

\begin{proposition}\label{prop:trivial-p1-tree}
The analytification of $\mathbf P^1_k$ is obtained by gluing, at their
$t=0$ endpoints, one closed ray
$\{v_{a,t}:0\leq t\leq+\infty\}$ for each closed point
$a\in\mathbf P^1(k)$.  Positive real scaling acts by
\[
                         \lambda\cdot v_{a,t}=v_{a,\lambda t}.
\]
The common vertex is the trivial valuation and the endpoint of each ray is a
proper-support semivaluation.
\end{proposition}

\begin{proof}
A valuation of the function field of a smooth curve centered at a closed
point is a nonnegative multiple of the order of vanishing at that point.
A semivaluation with nonzero support factors through the residue field of a
closed point and is therefore the displayed endpoint.  Pointwise convergence
on regular functions gives the asserted topology.  Along a fixed branch it
is the usual topology of $[0,+\infty]$.  A basic neighborhood of the common
vertex contains every branch except possibly finitely many, and contains an
initial interval on each exceptional branch: indeed, a rational function has
a zero or pole at only finitely many closed points.  As $t\to\infty$, a
function vanishing at $a$ acquires value $+\infty$, while a unit retains value
$0$.  Thus the gluing is understood with this compact, evaluation topology,
not with the noncompact weak-wedge topology on an infinite family of rays.
\end{proof}

The cone structure is not optional.  Over a nontrivially valued field the
condition $v(\pi)=1$ cuts out a preferred slice; over a trivially valued field
there is no such element.  Homogeneous functions satisfy
\[
                              \varphi(\lambda v)=\lambda\varphi(v).
\]

The next proposition makes this precise for the curve and shows exactly what
is lost by attempting to normalize.

\begin{proposition}
\label{prop:ch08-trivial-curve-scaling}
Let $X$ be a smooth connected projective curve over $k$ and let
$X^{\An}$ carry the trivial absolute value on $k$.
\begin{enumerate}
\item The map $(t,v)\mapsto tv$ is a continuous action of
$\mathbf R_{>0}$ on $X^{\An}$ whose fixed points are exactly the trivial
valuation $v_\triv$ of $k(X)$ and the proper-support semivaluations
$v_{a,\infty}$, $a\in X(k)$.
\item A point $v_{a,t}$ with $0<t<\infty$ is divisorial---that is, of the
form $c\,\ord_E$ for a divisorial valuation $\ord_E$ and $c\in\mathbf Q_{>0}$
---if and only if $t\in\mathbf Q_{>0}$.  Divisorial points are dense in
$X^{\An}$.
\item The subset $\Lambda=\{\ord_a:a\in X(k)\}$, the normalized link of
\cref{def:normalized-link} for the log-discrepancy function
$A(v_{a,t})=t$, is discrete in $X^{\An}$.  It is therefore compact if and
only if $X(k)$ is finite, which never happens for $k$ algebraically closed.
\end{enumerate}
\end{proposition}

\begin{proof}
(1)  Continuity of the action is \eqref{eq:valuation-scaling}.  By
\cref{prop:trivial-p1-tree} and its evident analogue for an arbitrary smooth
curve---every valuation of $k(X)$ trivial on $k$ and centred at a closed
point $a$ is a positive multiple of $\ord_a$, and a semivaluation with
nonzero support factors through a residue field---every point is of the form
$v_{a,t}$ with $t\in[0,+\infty]$, and $s\cdot v_{a,t}=v_{a,st}$ with
$s\cdot\infty=\infty$ and $s\cdot0=0$.  Hence the fixed points are those with
$t\in\{0,\infty\}$, that is, $v_\triv$ and the endpoints.

(2)  On a smooth curve the only divisorial valuations of $k(X)$ are the
$\ord_a$, $a\in X(k)$, up to positive rational multiples, because a prime
divisor on a birational model of a curve is a closed point of the curve
itself.  So $v_{a,t}=t\,\ord_a$ is divisorial precisely when
$t\in\mathbf Q_{>0}$.  Density holds along each ray, since
$\mathbf Q_{>0}$ is dense in $(0,\infty)$ and $t\mapsto v_{a,t}$ is
continuous with $v_{a,0}=v_\triv$ and
$\lim_{t\to\infty}v_{a,t}=v_{a,\infty}$, and every point lies on such a ray.

(3)  Fix $a\in X(k)$, choose $c\in X(k)$ with $c\neq a$, and put
$U=X\setminus\{c\}$, an affine open subset.  Every point of $X^{\An}$ other
than $v_{c,\infty}$ has its support and its centre in $U$, so
$U^{\An}=X^{\An}\setminus\{v_{c,\infty}\}$, an open subset containing all of
$\Lambda$.  Let $f_1,\ldots,f_N\in\cO_X(U)$ generate the ideal sheaf
$\mathfrak m_a$ of $a$ on $U$.  For $v\in U^{\An}$ the local generators at
$c(v)\in U$ may be taken among the $f_\ell$, so
\[
 \theta(v):=v(\mathfrak m_a)=\min_{1\leq\ell\leq N}v(f_\ell) ,
\]
which is continuous on $U^{\An}$ with values in $[0,+\infty]$, because each
$v\mapsto v(f_\ell)=-\log|f_\ell(v)|$ is continuous by the definition of the
topology of $U^{\An}$.  Now $\theta(v_{b,t})=t$ if $b=a$ and $\theta=0$
otherwise, since $\mathfrak m_a$ is the unit ideal at every $b\neq a$.
Therefore
\[
 U_a=\bigl\{v\in U^{\An}:\ \tfrac12<\theta(v)<\tfrac32\bigr\}
\]
is an open subset of $X^{\An}$ with $\Lambda\cap U_a=\{\ord_a\}$, and
$\Lambda$ is discrete.  A discrete subset of a Hausdorff space is compact if
and only if it is finite, and $X(k)$ is infinite because $k$ is
algebraically closed and $X$ is a curve.
\end{proof}

\begin{remark}[Why the theory is homogeneous]
\Cref{prop:ch08-trivial-curve-scaling}\,(3) is the structural reason that
\cref{part:bj} works on the whole cone $X^{\An}$ and with homogeneous
functions, rather than on a compact link with ordinary functions as in
\cref{part:bfj}.  There is no compact slice transverse to the scaling action.
What replaces it is the pair (homogeneous function on the cone, bounded
function on the compact space $X^{\An}$), related by the Legendre
correspondence of \cref{thm:homogeneous-legendre-decomposition}.
\end{remark}

\subsection{Homogeneous potentials of ideals}
\label{sec:ch08-trivial-potentials}
\index{potential!homogeneous}

Let $X$ be a \emph{proper} variety over the trivially valued $k$, so that
every $v\in X^{\An}$ has a centre $c(v)\in X$
(\cref{prop:centers-proper-separated}), and let
$\mathfrak a\subseteq\cO_X$ be a nonzero coherent ideal.  Following
\eqref{eq:ideal-evaluation-triv}, $v$ evaluates on
$\mathfrak a$ by
\[
 v(\mathfrak a)=\min\{v(f):f\in\mathfrak a_{c(v)}\}\in[0,+\infty] ,
\]
the minimum being attained on a finite generating set, and we set
\begin{equation}
\label{eq:ch08-ideal-potential}
 \varphi_{\mathfrak a}(v)=-v(\mathfrak a)\in[-\infty,0] .
\end{equation}

\begin{proposition}
\label{prop:ch08-ideal-potential}
With the notation above,
\begin{enumerate}
\item $\varphi_{\mathfrak a}$ is homogeneous:
$\varphi_{\mathfrak a}(tv)=t\varphi_{\mathfrak a}(v)$ for $t>0$;
\item $\varphi_{\mathfrak a\mathfrak b}
      =\varphi_{\mathfrak a}+\varphi_{\mathfrak b}$ and
      $\varphi_{\mathfrak a+\mathfrak b}
      =\max\{\varphi_{\mathfrak a},\varphi_{\mathfrak b}\}$;
\item $\varphi_{\mathfrak a}\leq0$, with equality at $v$ exactly when
$\mathfrak a$ is trivial at the centre of $v$;
\item if $X$ is a smooth projective curve and $c\in X(k)$, then
$\varphi_{\mathfrak a}$ is continuous, with values in $[-\infty,0]$, on the
open set $X^{\An}\setminus\{v_{c,\infty}\}$; in particular
$\varphi_{\mathfrak a}$ is continuous at every point of $X^{\An}$ other than
the proper-support endpoints.
\end{enumerate}
\end{proposition}

\begin{proof}
(1)  Immediate from $(tv)(f)=t\,v(f)$ and $t\cdot(+\infty)=+\infty$.

(2)  If $\mathfrak a$ is generated by $f_i$ and $\mathfrak b$ by $g_j$ near
the centre, then $\mathfrak a\mathfrak b$ is generated by the $f_ig_j$ and
$\min_{i,j}\{v(f_i)+v(g_j)\}=\min_iv(f_i)+\min_jv(g_j)$; likewise
$\mathfrak a+\mathfrak b$ is generated by the $f_i$ and $g_j$ together, and
the minimum of the union of two finite sets is the minimum of the two
minima.  Negating gives the stated identities; these are the tropical
identities \eqref{eq:ideal-tropical-identities}.

(3)  $v\geq0$ on the local ring at the centre, so $v(\mathfrak a)\geq0$;
it vanishes exactly when some generator is a unit at the centre, that is,
when $\mathfrak a_{c(v)}=\cO_{X,c(v)}$.

(4)  This is the computation carried out in the proof of
\cref{prop:ch08-trivial-curve-scaling}\,(3) for $\mathfrak a=\mathfrak m_a$,
verbatim for a general $\mathfrak a$: with $U=X\setminus\{c\}$ affine and
$f_1,\ldots,f_N\in\cO_X(U)$ generating $\mathfrak a|_U$, every point of
$U^{\An}=X^{\An}\setminus\{v_{c,\infty}\}$ has its centre in $U$, so
$v(\mathfrak a)=\min_\ell v(f_\ell)$ there; each $v\mapsto v(f_\ell)$ is
continuous into $[0,+\infty]$, and a minimum of finitely many continuous
functions is continuous.  Since $c$ was arbitrary, the last assertion
follows by varying $c$.
\end{proof}

\begin{example}[A monomial ideal]
Let $X=\mathbf A^2_k$ and let $v_w$ be the monomial valuation with
$v_w(x)=w_1$, $v_w(y)=w_2$ for $w\in\mathbf R_{\geq0}^2$.  For
$\mathfrak a=(x^2,y^3)$,
\[
 v_w(\mathfrak a)=\min\{2w_1,3w_2\},
 \qquad
 \varphi_{\mathfrak a}(v_w)=-v_w(\mathfrak a)
 =\max\{-2w_1,-3w_2\}.
\]
Thus a coherent ideal gives a homogeneous piecewise affine function on the
valuation cone.  The wall $2w_1=3w_2$ is the tropical trace of the competition
between its two generators.
\end{example}

\subsection{The monomial cone of the affine plane, computed}
\label{sec:ch08-monomial-cone}
\index{monomial valuation!cone of}

We continue with $X=\mathbf A^2_k$ and record the structure of the monomial
part of its valuation space, together with the invariants that
\cref{part:bj} attaches to a homogeneous potential.

For $w=(w_1,w_2)\in\mathbf R^2_{\geq0}$, let $v_w$ be the monomial valuation
with $v_w(x)=w_1$, $v_w(y)=w_2$, so that
\begin{equation}
\label{eq:ch08-monomial-value}
 v_w\Bigl(\sum_{\beta\in\mathbf N^2}c_\beta x^{\beta_1}y^{\beta_2}\Bigr)
 =\min_{c_\beta\neq0}\langle w,\beta\rangle ,
\end{equation}
which is a valuation by \cref{lem:monomial-valuation} (or by the direct
graded-domain argument of \cref{ex:graded-gauss}).  The map $w\mapsto v_w$
identifies the closed quadrant $\mathbf R^2_{\geq0}$ with a closed subset of
$(\mathbf A^2_k)^{\An}$, and the scaling action becomes scalar
multiplication:
\[
 t\cdot v_w=v_{tw},\qquad t>0 .
\]
The vertex $w=0$ is the trivial valuation.  The two boundary rays are
$w=(s,0)$, the valuations $s\,\ord_{\{x=0\}}$, and $w=(0,s)$; the interior
rays are the genuinely two-dimensional monomial valuations.

\begin{proposition}
\label{prop:ch08-monomial-invariants}
Let $w\in\mathbf R^2_{>0}$.
\begin{enumerate}
\item In the classical valuation-theoretic convention, which permits
arbitrary positive real rescaling, $v_w$ is divisorial up to scale if and
only if $w\in\mathbf R_{>0}\cdot\mathbf Q^2$, equivalently
$w_1/w_2\in\mathbf Q$.  In the rational divisorial-point convention of
\Cref{prop:ch08-trivial-curve-scaling}, it is divisorial if and only if
$w\in\mathbf Q_{>0}^2$.  If $w_1/w_2\notin\mathbf Q$, it is quasimonomial
of rational rank two and is not divisorial in either sense.
\item The log discrepancy of $v_w$ relative to $\mathbf A^2$ is
$A(v_w)=w_1+w_2$, and the normalized link
$\{v_w:w_1+w_2=1\}$ is the closed segment joining
$\ord_{\{x=0\}}$ to $\ord_{\{y=0\}}$.
\item For a monomial ideal $\mathfrak a=(x^{p_1}y^{q_1},\ldots,
x^{p_N}y^{q_N})$,
\[
 \varphi_{\mathfrak a}(v_w)=-\min_{\ell}\{p_\ell w_1+q_\ell w_2\}
   =\max_{\ell}\langle-(p_\ell,q_\ell),w\rangle ,
\]
a homogeneous, convex, piecewise linear function of $w$.
\end{enumerate}
\end{proposition}

\begin{proof}
(1)  The value group of $v_w$ is $\mathbf Z w_1+\mathbf Z w_2$.  A divisorial
valuation has value group of rational rank one, so
$\mathbf Zw_1+\mathbf Zw_2$ must have rational rank one, which forces
$w_1/w_2\in\mathbf Q$; conversely, if $w=c(p,q)$ with $p,q$ coprime positive
integers and $c>0$, then $v_{(p,q)}$ is the divisorial valuation
$\ord_E$ attached to the exceptional divisor of the weighted blowup of
weights $(p,q)$, and $v_w=c\,\ord_E$.  This rescaling is rational exactly
when $c\in\mathbf Q_{>0}$, equivalently when
$w\in\mathbf Q_{>0}^2$.  The rational-rank computation for
monomial valuations is \cref{lem:monomial-invariants}; the constraint linking
rational rank and residue transcendence degree on a surface is Abhyankar's
inequality \cref{thm:abhyankar-inequality}.

(2)  On the weighted blowup with weights $(p,q)$, coprime, the exceptional
divisor $E$ satisfies $K_{\cX/\mathbf A^2}=(p+q-1)E$, so
\eqref{eq:appJ-div-discrepancy} gives
$A_{\mathbf A^2}(\ord_E)=1+(p+q-1)=p+q$.  Homogeneity
\eqref{eq:appJ-discrepancy-homogeneous} extends this to all
$w\in\mathbf R_{>0}\cdot(\mathbf Q^2\cap\mathbf R^2_{>0})$, and continuity to
the closed quadrant.  The description of the link is then immediate, the
endpoints being $w=(1,0)$ and $w=(0,1)$.

(3)  Apply \eqref{eq:ch08-monomial-value} to each generator and use
\cref{prop:ch08-ideal-potential}\,(2): the value of $\mathfrak a$ is the
minimum of the values of the generators, so its negative is the maximum of
the negatives.
\end{proof}

\begin{example}[The multiplicity of $(x^2,y^3)$, computed twice]
\label{ex:ch08-monomial-multiplicity}
\index{multiplicity!of a monomial ideal}
Take $\mathfrak a=(x^2,y^3)$ as in the previous example, so that
\[
 \varphi_{\mathfrak a}(v_w)=\max\{-2w_1,\,-3w_2\} ,
\]
with wall $2w_1=3w_2$, that is, $w_2=\tfrac23w_1$.  Two computations of the
same number illustrate the pattern that recurs in the toric case.

\emph{Algebraic side.}  The ideal $\mathfrak a$ is generated by a regular
sequence, so its Samuel multiplicity equals its colength:
\[
 e(\mathfrak a)=\dim_kk[x,y]/(x^2,y^3)=2\cdot3=6 .
\]

\emph{Convex side.}  The Newton polyhedron of $\mathfrak a$ is
$\Gamma_+=\operatorname{conv}\bigl(((2,0)+\mathbf R^2_{\geq0})\cup
((0,3)+\mathbf R^2_{\geq0})\bigr)$, whose boundary staircase in the quadrant
is the segment from $(2,0)$ to $(0,3)$; the complementary region is
\[
 \mathbf R^2_{\geq0}\setminus\Gamma_+
 =\{(a,b)\geq0:3a+2b<6\} ,
\]
the open triangle with vertices $(0,0),(2,0),(0,3)$, of area
$\tfrac12\cdot2\cdot3=3$.  Hence
\[
 2!\operatorname{vol}\bigl(\mathbf R^2_{\geq0}\setminus\Gamma_+\bigr)
 =2\cdot3=6=e(\mathfrak a) .
\]
The two answers agree.  Here they were both computed; in general the
identity $e(\mathfrak a)=n!\operatorname{vol}(\mathbf R^n_{\geq0}
\setminus\Gamma_+)$ for a monomial ideal of finite colength is the classical
Newton-polyhedron formula of Teissier, which we do not use elsewhere.  The factor $n!=2$ is the same factor that appears in
\eqref{eq:toric-ma-comparison} and in \eqref{eq:appJ-toric-degree}: a
Monge--Amp\`ere or multiplicity number is $n!$ times a Euclidean volume.  In
the language of \cref{part:bj}, $6$ is the mass of the homogeneous
Monge--Amp\`ere measure of $\varphi_{\mathfrak a}$ in the local setting of
\cref{chap:local-valuative-psh}.
\end{example}

\begin{example}[Truncation and boundedness]
\label{ex:ch08-truncation}
Homogeneous potentials are unbounded on every scaling orbit unless they
vanish, by \cref{def:homogeneous-function}.  For $\mathfrak a=(x^2,y^3)$ and
$c>0$, the truncation
\[
 \max\{\varphi_{\mathfrak a},-c\}(v_w)
 =\max\bigl\{-2w_1,\,-3w_2,\,-c\bigr\}
\]
is bounded and continuous.  It equals $\varphi_{\mathfrak a}$ precisely on
\[
                  \{\min(2w_1,3w_2)\le c\},
\]
and equals $-c$ when $2w_1\ge c$ and $3w_2\ge c$.  Letting
$c\to+\infty$ recovers $\varphi_{\mathfrak a}$ as a decreasing limit of
bounded functions.  This is the mechanism used in
\cref{ex:ch01-cone-ideal,ex:ch01-product-tc}: a test configuration produces a
truncated homogeneous potential, and the truncation level is the
normalization of the associated Fubini--Study function.
\end{example}

\subsection{Monge--Amp\`ere measures on the trivially valued line}
\label{sec:ch08-trivial-ma}

On a smooth projective curve $X$ over the trivially valued $k$, the
tree structure of \cref{prop:trivial-p1-tree} makes the Monge--Amp\`ere
operator of \cref{part:bj} completely explicit: by
\eqref{eq:bj-curve-laplacian}, for a line bundle $L$ of degree $d$ and a
piecewise affine $\phi$ on the tree,
\begin{equation}
\label{eq:ch08-trivial-curve-ma}
 \omega_\triv+\ddc\phi=d\,\delta_{v_\triv}+\Delta\phi ,
\end{equation}
where $\Delta$ is the sum-of-outgoing-slopes Laplacian of
\eqref{eq:graph-laplacian}.  The tree itself is not locally finite over an
algebraically closed field, but the functions considered below have nonzero
slopes on only finitely many rays, so every displayed slope sum is finite.
Two computations follow; both should be
compared with \cref{ex:ch01-product-tc}.

\begin{example}[One attracting fixed point]
\label{ex:ch08-trivial-one-ray}
Let $X=\mathbf P^1_k$, $L=\cO(1)$, $d=1$, and fix $c>0$.  Define $\phi$ to be
$0$ on every ray except the ray through $0=[1:0]$, on which
\[
 \phi(v_{0,t})=\max\{-t,-c\} .
\]
The function is continuous, and it has exactly two corners.  At $v_\triv$
the only ray along which $\phi$ is nonconstant is the ray through $0$, with
outgoing slope $-1$; all other outgoing slopes vanish; so
$(\Delta\phi)(\{v_\triv\})=-1$.  At $v_{0,c}$ the two outgoing slopes are
$+1$ towards the vertex and $0$ away from it, so
$(\Delta\phi)(\{v_{0,c}\})=+1$.  Hence
\[
 \omega_\triv+\ddc\phi
 =1\cdot\delta_{v_\triv}-\delta_{v_\triv}+\delta_{v_{0,c}}
 =\delta_{c\,\ord_0} ,
\]
a single atom of mass $1=\deg\cO(1)$ at the divisorial point $c\,\ord_0$,
provided $c\in\mathbf Q_{>0}$.  For $c=1$ this is exactly the measure
computed in \cref{ex:ch01-product-tc} for the running product test
configuration.
\end{example}

\begin{example}[Splitting the mass between two rays]
\label{ex:ch08-trivial-two-rays}
Keep $X=\mathbf P^1_k$ and $L=\cO(1)$, and now put, for $c>0$,
\[
 \phi(v_{0,t})=\max\{-\tfrac t2,-\tfrac c2\},\qquad
 \phi(v_{\infty,t})=\max\{-\tfrac t2,-\tfrac c2\},\qquad
 \phi\equiv0\ \text{on all other rays.}
\]
At $v_\triv$ there are now two rays with outgoing slope $-\tfrac12$, so
$(\Delta\phi)(\{v_\triv\})=-1$; at each of $v_{0,c}$ and $v_{\infty,c}$ the
mass is $+\tfrac12$.  Therefore
\[
 \omega_\triv+\ddc\phi
 =\bigl(1-1\bigr)\delta_{v_\triv}
   +\tfrac12\delta_{c\,\ord_0}+\tfrac12\delta_{c\,\ord_\infty} ,
\]
of total mass $1$, with no mass at the vertex.  Trying instead to use slope
$-1$ on both rays would give
$(\Delta\phi)(\{v_\triv\})=-2$ and a total mass of $-1$ at the vertex: the
resulting function is not $\omega_\triv$-psh.  The correct constraint is
\begin{equation}
\label{eq:ch08-trivial-slope-constraint}
 -\sum_{\text{rays }a}\ d_a\phi(v_\triv)\ \leq\ d ,
\end{equation}
where $d_a\phi(v_\triv)\leq0$ are the outgoing slopes at the vertex.  This
inequality is the trivially valued avatar of the nef-interval condition
\eqref{eq:running-nef-interval}.
\end{example}

\subsection{The Gauss extension and the two pictures}
\label{sec:ch08-gauss-extension}

The trivially valued cone and the nontrivially valued skeleton are not
unrelated.  Put $K=k((\varpi))$ with $v_K(\varpi)=1$.  The Gauss extension
\eqref{eq:gauss-extension-triv} sends a full-support $v\in X^{\An}$ to the
point $\sigma(v)\in X_K^{\An}$ defined by
\[
 \sigma(v)\Bigl(\sum_if_i\varpi^i\Bigr)=\min_i\{v(f_i)+i\} .
\]

\begin{example}[The Gauss extension of a ray]
\label{ex:ch08-gauss-ray}
Take $X=\mathbf P^1_k$ and $v=t\,\ord_0$ with $t>0$, in the coordinate $T$
with $\ord_0(T)=1$.  Then $\sigma(v)$ is the monomial point of
$(\mathbf P^1_K)^{\An}$ with
\[
 \sigma(v)(T)=t,\qquad \sigma(v)(\varpi)=1 ,
\]
that is, the Gauss point $\zeta_{0,|\varpi|^t}$ of the disc of radius
$|\varpi|^t$, at tropical coordinate $t$ in the notation of
\cref{sec:ch08-p1-tree}.  As $t$ ranges over $[0,\infty)$, the image is the
half-apartment $\{t\geq0\}$; the vertex $v_\triv$ goes to the Gauss point of
the unit disc.  Thus the ray of the trivially valued picture is literally the
segment computed in \cref{ex:ch08-p1-interval,ex:ch08-newton-cubic}, and the
scaling action becomes the action of rescaling the ground valuation.  The
homogeneous potential $\varphi_{\mathfrak m_0}(t\,\ord_0)=-t$ of
\cref{ex:ch01-cone-ideal} becomes the model function $-\varphi_{E}$ of
\cref{ex:ch01-blowup-model-function} on the corresponding segment.  The
systematic treatment of this correspondence is
\cref{chap:trivial-analytification}.
\end{example}

On the ray $v_{a,t}$, a homogeneous function is determined by its value at
$t=1$.  General psh functions in the trivially valued theory are not all
homogeneous; the homogeneous part records $b$-divisorial data, while bounded
translations and Legendre decompositions carry the remaining information.

\section{Toric skeleta and convex metric functions}
\label{sec:toric-convex}

Let $N\simeq\mathbf Z^n$ be a lattice, $M=N^\vee$, and
$\mathbb T=\operatorname{Spec}K[M]$.  The tropicalization map is normalized
by
\begin{equation}
\label{eq:toric-trop}
 \langle m,\operatorname{trop}(x)\rangle
 =-\log|\chi^m(x)|,
 \qquad m\in M.
\end{equation}
For $u\in N_{\mathbf R}$, the Gauss section $\iota:N_{\mathbf R}\to
\mathbb T^{\mathrm{an}}$ is
\begin{equation}
\label{eq:toric-gauss-section}
 \left|\sum_m a_m\chi^m\right|_{\iota(u)}
 =\max_m |a_m|\rme^{-\langle m,u\rangle}.
\end{equation}
Then $\operatorname{trop}\circ\iota$ is the identity.  The image is the
toric skeleton of the dense torus.  The next lemma proves this and fixes the
elementary facts used in every computation below.

\begin{lemma}
\label{lem:ch08-toric-gauss-section}
Let $K$ be complete and nontrivially valued.
\begin{enumerate}
\item For every $u\in N_{\mathbf R}$, formula
\eqref{eq:toric-gauss-section} defines a multiplicative seminorm on $K[M]$,
hence a point $\iota(u)\in\mathbb T^{\An}$, and
$\operatorname{trop}(\iota(u))=u$.
\item $\iota$ is a continuous section of $\operatorname{trop}$, and it is a
homeomorphism onto its image.
\item For $f=\sum_ma_m\chi^m\in K[M]$ one has
\begin{equation}
\label{eq:ch08-toric-trop-poly}
 -\log|f(\iota(u))|
   =\min_m\bigl\{v(a_m)+\langle m,u\rangle\bigr\} ,
\end{equation}
a concave piecewise affine function of $u$ with slopes in $M$.
\end{enumerate}
\end{lemma}

\begin{proof}
(1)  Write $|f|_u=\max_m|a_m|\rme^{-\langle m,u\rangle}$.  The ultrametric
inequality is clear.  For multiplicativity, put
$\Lambda=\{m:|a_m|\rme^{-\langle m,u\rangle}=|f|_u\}$ and define $\Lambda'$
similarly for $g=\sum_mb_m\chi^m$.  In the graded reduction attached to the
weight $m\mapsto v(a_m)+\langle m,u\rangle$, the initial forms of $f$ and
$g$ are nonzero elements of a domain, namely the group algebra of $M$ over
the graded residue field of $K$
(\cref{lem:graded-residue-field,ex:graded-gauss}), because the group algebra
of a torsion-free abelian group over a domain is a domain.  Hence the initial
form of $fg$ is the product of the two initial forms and is nonzero, so
$|fg|_u=|f|_u|g|_u$.  Applying the definition to $f=\chi^m$ gives
$|\chi^m(\iota(u))|=\rme^{-\langle m,u\rangle}$, which by
\eqref{eq:toric-trop} is exactly $\operatorname{trop}(\iota(u))=u$.

(2)  $\iota$ is continuous because each function
$u\mapsto|f(\iota(u))|$ is a maximum of finitely many continuous functions of
$u$, and $\operatorname{trop}$ is continuous by \eqref{eq:toric-trop}.  A
continuous section of a continuous map is a homeomorphism onto its image,
with inverse the restriction of $\operatorname{trop}$.

(3)  Take $-\log$ of the definition and use $v=-\log|\cdot|$; the minimum of
finitely many affine functions of $u$ with linear parts $\langle m,\cdot\rangle$
is concave and piecewise affine with slopes in $M$.
\end{proof}

For $N=\mathbf Z$, formula \eqref{eq:ch08-toric-trop-poly} is the Newton
polygon computation of \cref{ex:ch08-newton-cubic}; the general case is the
tropicalization of a Laurent polynomial, and its corner locus is the
tropical hypersurface of $f$.

Let $X_\Sigma$ be a proper toric variety and let $L$ be an ample toric line
bundle.  After choosing a toric section and a sign convention, $L$ is encoded
by a full-dimensional lattice polytope $P\subset M_{\mathbf R}$.  We use the
convex support function and its reflection
\[
 h_P(u)=\sup_{m\in P}\langle m,u\rangle,
 \qquad
 \check h_P(u):=h_P(-u)=\sup_{m\in P}\langle-m,u\rangle.
\]
The reflection is forced by our conventions: tropical coordinates are
$u=-\log|\chi|$, whereas metric weights are $-\log\|s\|$.  Thus the
recession function of a weight is $\check h_P$, not $h_P$.  For example,
$P=[0,d]$ gives $\check h_P(u)=\max\{0,-du\}$, in agreement with
\eqref{eq:p1-standard-weight}.

\begin{definition}
A \emph{toric convex metric function} for $(L,P)$ is a convex function
$F:N_{\mathbf R}\to\mathbf R$ such that $F-\check h_P$ is bounded.  It is a
\emph{model metric function} when it is rational piecewise affine and its
affine constants lie in the divisible hull of $v(K^\times)$, so that a
positive multiple is induced by a toric model.
\end{definition}

Changing the toric section adds an affine function to both $F$ and
$\check h_P$;
the metric and its Monge--Amp\`ere measure do not change.  Convexity is the
toric form of semipositivity, while the bounded-recession condition fixes the
line bundle.

\subsection{Alexandrov measure}

For a finite convex function $F$ on $N_{\mathbf R}$, its subgradient at $u$
is
\[
 \partial F(u)=
 \{m\in M_{\mathbf R}:F(w)\geq F(u)+\langle m,w-u\rangle
   \text{ for all }w\}.
\]
The real Alexandrov Monge--Amp\`ere measure is characterized by
\begin{equation}
\label{eq:alexandrov-ma}
                 \operatorname{MA}_{\mathbf R}(F)(E)
                 =\operatorname{vol}_M(\partial F(E))
\end{equation}
for Borel $E$, with the standard disjointification on overlaps of
subgradients.  Lebesgue measure is normalized so that a fundamental
parallelepiped of $M$ has volume $1$.

If $F$ is $C^2$ and strictly convex, then
\begin{equation}
\label{eq:smooth-real-ma}
             \operatorname{MA}_{\mathbf R}(F)
             =\det(D^2F)\,\rd u.
\end{equation}
If $F$ is piecewise affine, the measure is atomic and
\begin{equation}
\label{eq:pl-real-ma}
 \operatorname{MA}_{\mathbf R}(F)(\{u\})
 =\operatorname{vol}_M(\partial F(u)).
\end{equation}

The word ``disjointification'' in \eqref{eq:alexandrov-ma} deserves an
explanation, because it is the only point at which the definition could fail
to produce a measure.

\begin{lemma}
\label{lem:ch08-alexandrov-polyhedral}
Let $F=\max_{1\leq i\leq N}\ell_i$ be a finite maximum of affine functions
$\ell_i(u)=\langle m_i,u\rangle+c_i$ on $N_{\mathbf R}$, with distinct
$m_i\in M_{\mathbf R}$.  For $u\in N_{\mathbf R}$ put
$I(u)=\{i:\ell_i(u)=F(u)\}$.
\begin{enumerate}
\item $\partial F(u)=\operatorname{conv}\{m_i:i\in I(u)\}$.
\item If $u\neq u'$, then $\partial F(u)\cap\partial F(u')$ has empty
interior; hence $\operatorname{vol}_M(\partial F(u))>0$ for at most finitely
many $u$, and \eqref{eq:pl-real-ma} defines a finite atomic measure.
\item $\operatorname{vol}_M(\partial F(u))>0$ if and only if
$\{m_i:i\in I(u)\}$ affinely spans $M_{\mathbf R}$, that is, if and only if
$u$ is a vertex of the polyhedral decomposition of $N_{\mathbf R}$ into the
domains of linearity of $F$.
\end{enumerate}
\end{lemma}

\begin{proof}
(1)  If $m=\sum_{i\in I(u)}\lambda_im_i$ with $\lambda_i\geq0$ summing to
$1$, then for every $w$,
\[
 \begin{aligned}
  F(w)&\geq\sum_{i\in I(u)}\lambda_i\ell_i(w)\\
      &=\sum_{i\in I(u)}\lambda_i
        \bigl(\ell_i(u)+\langle m_i,w-u\rangle\bigr)\\
      &=F(u)+\langle m,w-u\rangle .
 \end{aligned}
\]
so $m\in\partial F(u)$.  Conversely, $\partial F(u)$ is compact and convex,
and
\[
 \partial F(u)\subseteq\operatorname{conv}\{m_1,\ldots,m_N\}.
\]
A supporting-hyperplane argument identifies its extreme points among the
$m_i$ with $i\in I(u)$.  Indeed, if $m\in\partial F(u)$ and
$\xi\in N_{\mathbf R}$, then
\[
 \langle m,\xi\rangle
 \leq\lim_{s\downarrow0}s^{-1}\bigl(F(u+s\xi)-F(u)\bigr)
 =\max_{i\in I(u)}\langle m_i,\xi\rangle.
\]
Now a compact convex set whose
support function is dominated by that of
$\operatorname{conv}\{m_i:i\in I(u)\}$ is contained in it
\cite[Section~1.7]{Sch14}.

(2)  Suppose $m$ lies in the interior of $\partial F(u)\cap\partial F(u')$.
Then $F(u')\geq F(u)+\langle m,u'-u\rangle$ and
$F(u)\geq F(u')+\langle m,u-u'\rangle$, so both are equalities, and
therefore $m$ is a subgradient at every point of the segment $[u,u']$.
Choose $\xi=u'-u\neq0$; since $m$ is interior to $\partial F(u)$, also
$m+\varepsilon\xi^\vee\in\partial F(u)$ for a small $\varepsilon>0$ and any
$\xi^\vee\in M_{\mathbf R}$ with $\langle\xi^\vee,\xi\rangle>0$, whence
$F(u')\geq F(u)+\langle m,\xi\rangle+\varepsilon\langle\xi^\vee,\xi\rangle
>F(u')$, a contradiction.  For the finiteness, the sets
$\operatorname{conv}\{m_i:i\in I\}$ for $I\subseteq\{1,\ldots,N\}$ are
finitely many, and $\partial F(u)$ is one of them.

(3)  A polytope in $M_{\mathbf R}\simeq\mathbf R^n$ has positive Lebesgue
measure exactly when it is $n$-dimensional, that is, when its vertices
affinely span.
\end{proof}

\begin{lemma}\label{lem:toric-real-total-mass}
If $F-\check h_P$ is bounded, then
\[
                 \operatorname{MA}_{\mathbf R}(F)(N_{\mathbf R})
                 =\operatorname{vol}_M(P).
\]
\end{lemma}

\begin{proof}
The bounded difference gives the same recession function for $F$ and
$\check h_P$.
Convex duality then shows that the closure of the total subgradient image of
$F$ is $-P$.  The boundary of $-P$ has Lebesgue measure zero and
$\operatorname{vol}_M(-P)=\operatorname{vol}_M(P)$, so
\eqref{eq:alexandrov-ma} gives the formula.
\end{proof}

The factor relating real and non-Archimedean mass is now forced by toric
intersection theory.

\begin{theorem}\label{thm:toric-ma-comparison}
Let $X_\Sigma$ be a proper $n$-dimensional toric variety over a complete
discretely valued field, let $L$ be ample with polytope $P$, and let a
semipositive toric $\mathbf Q$-model metric be represented by the model
metric function $F$.  Then its Monge--Amp\`ere measure is supported on the
toric skeleton and
\begin{equation}
\label{eq:toric-ma-comparison}
 \operatorname{trop}_*c_1(L,\|\cdot\|)^n
 =n!\operatorname{MA}_{\mathbf R}(F).
\end{equation}
In particular, both sides have mass
\[
                         n!\operatorname{vol}_M(P)=(L^n).
\]
\end{theorem}

\begin{proof}[Toric measure-comparison input]
This is \cite[Theorem~5.70]{BGPS14}.  That theorem identifies the pushforward
of the Chambert--Loir measure of a semipositive algebraic toric metric with
$n!$ times the real Monge--Amp\`ere measure of its toric metric function.
Their function is concave; applying their formula to the reflected negative
of that function gives the convex convention $F$ used here.  Their valuation
map is our $\operatorname{trop}$, and their lattice-normalized measure is
$\operatorname{MA}_{\mathbf R}$.  Thus their formula is exactly
\eqref{eq:toric-ma-comparison}.  Its mass is
$n!\operatorname{vol}_M(P)=(L^n)$ by toric intersection theory.
\end{proof}

\begin{corollary}\label{cor:toric-ma-continuous}
Suppose that semipositive toric model metrics represented by $F_j$ converge
uniformly to a continuous metric represented by $F$.  Then
\[
 \operatorname{trop}_*c_1(L,\|\cdot\|)^n
 =n!\operatorname{MA}_{\mathbf R}(F).
\]
\end{corollary}

\begin{proof}
With the toric section fixed, uniform convergence of metrics is uniform
convergence of $F_j-F$.  Alexandrov measures of finite convex functions
converge weakly under locally uniform convergence.  Chambert--Loir measures
of semipositive metrics converge weakly under uniform convergence.  Passing
to the limit in \eqref{eq:toric-ma-comparison} proves the claim.  The
systematic invariant theory, including the global
tropical/non-Archimedean comparison beyond this approximation statement, is
developed in \cite{BGJKK25}.
\end{proof}

\begin{remark}[Exact scope of the skeletal comparison]
Vilsmeier's comparison \cite[Corollary~5.10]{Vil20} is local and has a
different hypothesis.  It starts with a strongly nondegenerate polystable
formal model, an $n$-dimensional open face $\tau$ of its skeleton with
corresponding special-fibre stratum $S$, and a convex function $h$ on
$\tau$.  On the inverse image of $\tau$, the metric
$\|1\|=\rme^{-h\circ p}$ on the trivial bundle satisfies
\[
 c_1(\mathcal O,\|\cdot\|)^n
 =\deg(S)\,n!\operatorname{MA}_{\mathbf R}(h).
\]
The factor $\deg(S)$ belongs to the skeletal calibration.  We use this
result later on polystable faces; it is not being invoked here as an
unrestricted global toric comparison theorem.
\end{remark}

\begin{example}[The support metric]
For $F=\check h_P$, one has $\partial F(0)=-P$.  Away from
lower-dimensional cones, the subgradient is contained in a proper face of
$-P$.  Hence
\[
              \operatorname{MA}_{\mathbf R}(\check h_P)
              =\operatorname{vol}(P)\delta_0
\]
and
\[
              c_1(L,\|\cdot\|_{\mathrm{can}})^n
              =n!\operatorname{vol}(P)\delta_{\iota(0)}.
\]
For $X=\mathbf P^1$ and $P=[0,d]$, this recovers
$d\delta_{\zeta_{0,1}}$.
\end{example}

\begin{example}[A smooth approximation on the toric line]
For $P=[0,d]$ and $\varepsilon>0$, set
\[
                  F_\varepsilon(u)
                  =\varepsilon\log(1+\rme^{-du/\varepsilon}).
\]
Then $F_\varepsilon-\max\{0,-du\}$ is bounded by
$\varepsilon\log2$, and
\[
 F_\varepsilon''(u)
 =\frac{d^2}{\varepsilon}
   \frac{\rme^{-du/\varepsilon}}{(1+\rme^{-du/\varepsilon})^2}.
\]
The integral of this density is $d$, since
$F_\varepsilon'$ increases from $-d$ to $0$.  As
$\varepsilon\downarrow0$, the measures converge weakly to $d\delta_0$.
This is the real convex model of concentration of a smooth approximation
onto a model skeleton.
\end{example}

\begin{example}[The cube]
Let $P=[0,1]^n$ and
\[
             F_\varepsilon(u)
             =\sum_{j=1}^n\varepsilon\log(1+\rme^{-u_j/\varepsilon}).
\]
The Hessian is diagonal, so
\[
 \det D^2F_\varepsilon
 =\prod_{j=1}^n
   \frac{1}{\varepsilon}
   \frac{\rme^{-u_j/\varepsilon}}{(1+\rme^{-u_j/\varepsilon})^2}.
\]
Its integral is $1=\operatorname{vol}(P)$, and it converges to $\delta_0$.
The non-Archimedean top wedge has mass $n!$, the degree of the corresponding
polarization.
\end{example}

\subsection{Worked examples in coordinates}
\label{sec:ch08-toric-worked}

We now carry out four computations completely.  In each of them the polytope
$P$, the reflected support function $\check h_P$, an explicit piecewise
affine perturbation, its subgradients, the resulting atoms, and the total
mass are written down, and the mass is checked against $(L^n)$.  The
governing identities are \eqref{eq:toric-ma-comparison} and
\[
 (L_P^n)=n!\operatorname{vol}_M(P),
\]
which is \eqref{eq:appJ-toric-degree}.

\begin{example}[Projective space]
\label{ex:ch08-toric-pn}
\leavevmode\par\noindent
Let
\[
 N=\mathbf Z^n,\qquad M=N^\vee,
 \qquad
 P=\operatorname{conv}\{0,e_1^\vee,\ldots,e_n^\vee\}.
\]
Thus $P$ is the standard simplex, the polytope of
$(\mathbf P^n,\cO(1))$.  Then
$h_P(u)=\max\{0,u_1,\ldots,u_n\}$ and
\[
 \check h_P(u)=\max\{0,-u_1,\ldots,-u_n\} .
\]
Its only subgradient of positive volume sits at $u=0$, where
$\partial\check h_P(0)=-P$, so by \eqref{eq:pl-real-ma}
\[
 \MA_{\mathbf R}(\check h_P)=\operatorname{vol}_M(P)\,\delta_0
 =\frac{1}{n!}\delta_0 ,
\]
since the standard simplex has lattice volume $1/n!$.  Multiplying by $n!$,
\[
 \operatorname{trop}_*c_1\bigl(\cO(1),\|\cdot\|_{\mathrm{std}}\bigr)^n
 =\delta_{\iota(0)} ,
\]
a Dirac mass of total mass $1=(\cO(1)^n)$ at the Gauss point.  This is
exactly the good-reduction computation of \cref{chap:metrics}, obtained there
from the fact that the special fibre of $\mathbf P^n_R$ is irreducible,
reduced, and of degree one.  The reflected support function is
\eqref{eq:p1-standard-weight} for $n=1$.
\end{example}

\begin{example}[$\mathbf P^1\times\mathbf P^1$]
\label{ex:ch08-toric-p1p1}
\index{toric geometry!quadric surface}
Let $n=2$ and $P=[0,a]\times[0,b]$ with $a,b\in\mathbf Z_{>0}$, the polytope
of $\bigl(\mathbf P^1\times\mathbf P^1,\ \cO(a,b)\bigr)$.  Then
\[
 \check h_P(u)=\max\{0,-au_1\}+\max\{0,-bu_2\},
 \qquad
 \operatorname{vol}_M(P)=ab ,
\]
and
\[
 2!\operatorname{vol}_M(P)=2ab=(aH_1+bH_2)^2 ,
\]
the last equality because $H_1^2=H_2^2=0$ and $(H_1\cdot H_2)=1$ on the
quadric surface.  The canonical model metric therefore has
$\MA_{\mathbf R}=ab\,\delta_0$ and non-Archimedean measure
$2ab\,\delta_{\iota(0)}$.

Now take $a=b=1$, so $P=[0,1]^2$ and $(L^2)=2$, and perturb.  For
$\lambda\geq0$ set
\[
 F_\lambda(u)=\max\{\,\ell_0,\ell_1,\ell_2,\ell_3\},
 \qquad
 \begin{aligned}
  \ell_0&=0, &\ell_1&=-u_1,\\
  \ell_2&=-u_2, &\ell_3&=-u_1-u_2+\lambda .
 \end{aligned}
\]
The four linear parts are $(0,0),(-1,0),(0,-1),(-1,-1)$, the vertices of
$-P$, so the recession function of $F_\lambda$ is
$\max\{0,-u_1,-u_2,-u_1-u_2\}=\check h_P(u)$ and
$0\leq F_\lambda-\check h_P\leq\lambda$: thus $F_\lambda$ is a toric convex
metric function for $\cO(1,1)$, and a model metric function when
$\lambda\in v(K^\times)$.

By \cref{lem:ch08-alexandrov-polyhedral}\,(3) the atoms sit at the points
where at least three of the $\ell_i$ are simultaneously maximal.  There are
four triples to test, and we test each:
\[
\begin{array}{lll}
 \ell_0=\ell_1=\ell_2:& u=(0,0),
  &\ell_3(0,0)=\lambda>0=\ell_0:\ \text{not maximal};\\
 \ell_0=\ell_1=\ell_3:& u=(0,\lambda),
  &\ell_2(0,\lambda)=-\lambda<0:\ \text{maximal};\\
 \ell_0=\ell_2=\ell_3:& u=(\lambda,0),
  &\ell_1(\lambda,0)=-\lambda<0:\ \text{maximal};\\
 \ell_1=\ell_2=\ell_3:& u=(\lambda,\lambda),
  &\ell_0=0>-\lambda:\ \text{not maximal}.
\end{array}
\]
(For $\lambda>0$; the degenerate case $\lambda=0$ gives back
$\check h_P$.)  Hence
\[
 \begin{aligned}
 \partial F_\lambda(0,\lambda)
   &=\operatorname{conv}\{(0,0),(-1,0),(-1,-1)\},\\
 \partial F_\lambda(\lambda,0)
   &=\operatorname{conv}\{(0,0),(0,-1),(-1,-1)\},
 \end{aligned}
\]
two triangles of lattice area $\tfrac12$ each.  Therefore
\[
 \MA_{\mathbf R}(F_\lambda)
 =\tfrac12\delta_{(0,\lambda)}+\tfrac12\delta_{(\lambda,0)},
 \qquad
 \MA_{\mathbf R}(F_\lambda)(N_{\mathbf R})=1=\operatorname{vol}_M(P) ,
\]
in agreement with \cref{lem:toric-real-total-mass}, and by
\eqref{eq:toric-ma-comparison}
\[
 \operatorname{trop}_*c_1\bigl(\cO(1,1),\|\cdot\|_{F_\lambda}\bigr)^2
 =2!\left(\tfrac12\delta_{(0,\lambda)}+\tfrac12\delta_{(\lambda,0)}\right)
 =\delta_{\iota(0,\lambda)}+\delta_{\iota(\lambda,0)} ,
\]
of total mass $2=(L^2)$.  The two triangles tile $-P$, which is the
geometric form of mass conservation: the subgradient polytopes of a convex
function with recession $\check h_P$ tile $-P$ up to sets of measure zero.
\end{example}

\begin{example}[A Hirzebruch surface]
\label{ex:ch08-toric-hirzebruch}
\index{Hirzebruch surface}
Let $n=2$ and
\[
 \begin{aligned}
 P&=\{m\in M_{\mathbf R}:m_1\geq0,\ m_2\geq0,\
       m_2\leq1,\ m_1+m_2\leq2\}\\
  &=\operatorname{conv}\{(0,0),(2,0),(1,1),(0,1)\} .
 \end{aligned}
\]
Its inner normal fan has rays $(1,0),(0,1),(-1,-1),(0,-1)$; since
$(0,-1)=(1,0)+(-1,-1)$, this fan is the star subdivision of the fan of
$\mathbf P^2$ at one maximal cone, so $X_\Sigma$ is the Hirzebruch surface
$\mathbf F_1=\operatorname{Bl}_p\mathbf P^2$, and $P$ is the polytope of
$L=2H-E$.  The polytope is a trapezoid with parallel sides of lengths $2$
and $1$ at distance $1$, so
\[
 \operatorname{vol}_M(P)=\frac{2+1}{2}=\frac32 ,
 \qquad
 2!\operatorname{vol}_M(P)=3=(2H-E)^2=4-1 ,
\]
a first mass check.  The reflected support function, computed on the four
vertices of $P$, is
\[
 \check h_P(u)=\max\{0,\,-2u_1,\,-u_1-u_2,\,-u_2\} .
\]

Perturb it by lowering the third and fourth affine pieces.  Set
\[
 F=\max\{\ell_0,\ell_1,\ell_2,\ell_3\},\qquad
 \begin{aligned}
  \ell_0&=0,           &\ell_1&=-2u_1+2,\\
  \ell_2&=-u_2+1,      &\ell_3&=-u_1-u_2+1 ,
 \end{aligned}
\]
whose linear parts $(0,0),(-2,0),(0,-1),(-1,-1)$ are the vertices of $-P$;
hence $F$ has recession function $\check h_P$ and differs from it by a
bounded amount, so it is a model metric function for $L$ as soon as
$1,2\in v(K^\times)$.  Testing the four triples:
\[
\begin{array}{lll}
 \ell_0=\ell_1=\ell_2:& u=(1,1),
  &\ell_3(1,1)=-1<0:\ \text{maximal};\\
 \ell_0=\ell_1=\ell_3:& u=(1,0),
  &\ell_2(1,0)=1>0:\ \text{not maximal};\\
 \ell_0=\ell_2=\ell_3:& u=(0,1),
  &\ell_1(0,1)=2>0:\ \text{not maximal};\\
 \ell_1=\ell_2=\ell_3:& u=(0,-1),
  &\ell_0=0<2=\ell_1(0,-1):\ \text{maximal}.
\end{array}
\]
So there are exactly two atoms, at $u=(1,1)$ and $u=(0,-1)$, with
\[
 \begin{aligned}
 \partial F(1,1)
   &=\operatorname{conv}\{(0,0),(-2,0),(0,-1)\},\\
 \partial F(0,-1)
   &=\operatorname{conv}\{(-2,0),(0,-1),(-1,-1)\} .
 \end{aligned}
\]
Their lattice areas are
\[
 \operatorname{vol}_M\bigl(\partial F(1,1)\bigr)
 =\tfrac12\bigl|\det\begin{psmallmatrix}-2&0\\0&-1\end{psmallmatrix}\bigr|
 =1,
 \qquad
 \operatorname{vol}_M\bigl(\partial F(0,-1)\bigr)
 =\tfrac12\bigl|\det\begin{psmallmatrix}2&-1\\1&-1\end{psmallmatrix}\bigr|
 =\tfrac12 ,
\]
the second determinant being formed from the edge vectors
$(0,-1)-(-2,0)=(2,-1)$ and $(-1,-1)-(-2,0)=(1,-1)$.  Hence
\[
 \MA_{\mathbf R}(F)=\delta_{(1,1)}+\tfrac12\delta_{(0,-1)},
 \qquad
 \text{total mass }\ 1+\tfrac12=\tfrac32=\operatorname{vol}_M(P) ,
\]
and by \eqref{eq:toric-ma-comparison}
\[
 \operatorname{trop}_*c_1(L,\|\cdot\|_F)^2
 =2\,\delta_{\iota(1,1)}+1\cdot\delta_{\iota(0,-1)} ,
 \qquad
 \text{total mass }\ 3=(L^2) .
\]
The corner locus of $F$ is the graph with the two vertices $(1,1)$ and
$(0,-1)$, joined by the segment on which $\ell_1=\ell_2$, namely
$u_2=2u_1-1$, together with four unbounded edges; it is the tropical
skeleton of the corresponding toric model over $K^\circ$, whose special
fibre has two irreducible surfaces of residual degrees $2$ and $1$.
\end{example}

\begin{example}[A weighted projective plane and rational masses]
\label{ex:ch08-toric-weighted}
\index{weighted projective space}
Let $X=\mathbf P(1,1,2)$, whose fan has rays $(1,0),(0,1),(-1,-2)$ in
$N=\mathbf Z^2$.  The ample generator $\cO_X(1)$ of the class group is a
rank-one reflexive sheaf representing an ample $\mathbf Q$-Cartier class,
not a line bundle.  Its polytope is
\[
 P_1=\operatorname{conv}\{(0,0),(1,0),(0,\tfrac12)\},
 \qquad
 \operatorname{vol}_M(P_1)=\tfrac14 ,
\]
so that the self-intersection of this $\mathbf Q$-Cartier class is
$(\cO_X(1)^2)=2!\cdot\tfrac14=\tfrac12$, a rational number, as it
must be on a singular toric surface.  The Cartier bundle
$L=\cO_X(2)$ has polytope $P=2P_1
=\operatorname{conv}\{(0,0),(2,0),(0,1)\}$, of lattice volume $1$, and
\[
 (L^2)=2!\cdot1=2 ,
\]
consistent with $(\cO_X(2)^2)=4\cdot\tfrac12=2$.  Its reflected support
function is
\[
 \check h_P(u)=\max\{0,\,-2u_1,\,-u_2\} ,
\]
with $\MA_{\mathbf R}(\check h_P)=\delta_0$ and non-Archimedean measure
$2\delta_{\iota(0)}$, of mass $2=(L^2)$.

The interest of this example is the failure of $\cO_X(1)$ to be Cartier.  Its
reflected support function $\check h_{P_1}=\tfrac12\check h_P$ is not
integral: at $u=(0,-1)$ it takes the value $\tfrac12$.  Correspondingly, the
model function is only a $\mathbf Q$-model function in the sense of
\cref{sec:metrics-model}, and the Monge--Amp\`ere mass
$2!\operatorname{vol}_M(P_1)=\tfrac12$ is not an integer.  Every statement in
this book about integral model line bundles is therefore stated for Cartier
divisors, with rational multiples allowed explicitly and never silently.
\end{example}

\begin{remark}[Which toric model realizes a given $F$]
\label{rem:ch08-toric-model}
A rational polyhedral convex $F$ with recession $\check h_P$ and constants in
$v(K^\times)$ determines a fan in $N_{\mathbf R}\times\mathbf R_{\geq0}$,
namely the normal fan of the epigraph
$\{(m,s):\langle m,u\rangle\leq F(u)+s\ \forall u\}
 =\{(m,s):s\geq F^{*}(m)\}$ of the Legendre transform
$F^{*}(m)=\sup_u(\langle m,u\rangle-F(u))$, whose domain is $-P$; the sign
matters, since with $-F$ in place of $F$ the set is empty whenever $P$ is
full-dimensional, positive homogeneity of $\check h_P$ forcing
$\langle m,u\rangle\leq-\check h_P(u)$ for all $u$ and hence
$\sup_{m'\in P}\langle m',u\rangle\leq\inf_{m'\in P}\langle m',u\rangle$.
The normal fan of that epigraph is a rational polyhedral fan in
$N_{\mathbf R}\times\mathbf R_{\geq0}$, and hence defines a toric
scheme over $K^\circ$ whose generic fibre is $X_\Sigma$ and whose special
fibre is the union of the toric surfaces dual to the vertices of the corner
locus of $F$.  In \cref{ex:ch08-toric-hirzebruch} the two vertices $(1,1)$
and $(0,-1)$ give the two components of the special fibre, and the residual
degrees $2$ and $1$ computed there are their self-intersection data.  The
construction is the toric case of \cref{thm:skeleton-subdivision}; a
systematic treatment is in \cite{BGPS14,CLS11}.
\end{remark}

\section{The running semistable degeneration}
\label{sec:running-semistable}

Let $R=K^\circ$ be a discrete valuation ring with $v(\pi)=1$.  The local
semistable equation
\begin{equation}
\label{eq:running-node}
                              xy=\pi
\end{equation}
has analytic generic fibre the annulus
$|\pi|\leq|x|\leq1$.  Put
\[
                   s=v(x)=-\log|x|.
\]
Then $0\leq s\leq1$ and $v(y)=1-s$.  The Gauss points with these values form
the skeleton $\Gamma=[0,1]$.  Retraction sends a point of the annulus to the
Gauss point having the same value of $|x|$.

The endpoints $v_0,v_1$ correspond to the two irreducible components of the
special fibre.  More generally, an equation $xy=a$ contributes an edge of
length $v(a)$.  Replacing a semistable model by a blowup subdivides the edge
but does not alter its path metric or the associated piecewise affine
function.

\begin{proposition}
\label{prop:ch08-annulus-skeleton}
Let $a\in K^\circ$ with $v(a)=\ell>0$ and let
\[
 A(\ell)=\{z\in\mathbf A^{1,\An}_K:0\leq v(z)\leq\ell\}
\]
be the corresponding closed annulus, the analytic generic fibre of
$\Spf R\{x,y\}/(xy-a)$ with $x=z$ and $y=a/z$.  Then:
\begin{enumerate}
\item the map $s\mapsto\zeta_{0,\rme^{-s}}$ identifies $[0,\ell]$ with a
closed subset $\Gamma_{A}\subset A(\ell)$, and the retraction
$r(z)=\zeta_{0,|z|}$ is a continuous idempotent onto it;
\item for $f$ in the annulus algebra, $|f(z)|\leq|f(r(z))|$, and
$s\mapsto-\log|f(\zeta_{0,\rme^{-s}})|$ is concave and piecewise affine with
integral slopes;
\item $\Gamma_A$ carries the path metric of \cref{def:ch08-path-metric}, in
which it has length $\ell=v(a)$; in particular the edge length is the
valuation of the deformation parameter, and is independent of the model.
\end{enumerate}
\end{proposition}

\begin{proof}
(1)  The formula $\zeta_{0,\rme^{-s}}$ makes sense for $s\in[0,\ell]$ and
lands in $A(\ell)$ since $v(z)=s$ at that point.  Continuity and idempotence
of $r$ are \cref{prop:skeleton-retraction} for the model
$\Spf R\{x,y\}/(xy-a)$, or directly: $|z|$ is a continuous function on
$A(\ell)$, and $r$ is the composition of $z\mapsto|z|$ with the continuous
map $\rho\mapsto\zeta_{0,\rho}$.

(2)  A function on the annulus is a convergent Laurent series
$f=\sum_{i\in\mathbf Z}c_iz^i$ with $|c_i|\rho^i\to0$ as $|i|\to\infty$ for
each $\rho\in[\rme^{-\ell},1]$.  The ultrametric inequality gives
$|f(z)|\leq\max_i|c_i||z|^i=|f(r(z))|$, and
\[
 -\log|f(\zeta_{0,\rme^{-s}})|=\min_i\{v(c_i)+is\}
\]
is a minimum of affine functions of $s$ with integral slopes, hence concave
and piecewise affine on the compact interval $[0,\ell]$; the decay condition
makes the minimum attained.

(3)  By \cref{def:ch08-path-metric} the point $\zeta_{0,\rme^{-s}}$ has
tropical coordinate $t=s$, so the length of $\Gamma_A$ is $\ell$.  The
number $\ell=v(a)$ is intrinsic: it is $v$ of the parameter, and a
different presentation $x'y'=a'$ of the same annulus has $v(a')=v(a)$ because
the annulus determines its two boundary Gauss points and hence their
distance.
\end{proof}

\begin{example}[The $A_{n-1}$ singularity and the length-$n$ edge]
\label{ex:ch08-An-edge}
Take $a=\pi^n$, so that the local model is
$\mathcal U_n=\Spec R[x,y]/(xy-\pi^n)$ and the edge has length $n$.  For $n\geq2$
the total space is normal but singular at the origin of the special fibre;
\cref{ex:snc-chain-models} resolves it by $\lfloor n/2\rfloor$ point blowups
and produces a chain
\[
 E_0,E_1,\ldots,E_n
\]
of $n+1$ components with all multiplicities one, the divisorial point of
$E_i$ being the monomial valuation with $v(x)=i$, $v(y)=n-i$.  The dual
complex is the segment $[0,n]$ subdivided at the integers, in agreement with
\cref{prop:ch08-annulus-skeleton}\,(3): resolving the singularity subdivides
the edge but does not change its length.
\end{example}

\begin{remark}[Edge lengths and base change]
\label{rem:ch08-base-change-lengths}
Let $K'=K(\pi^{1/m})$ with $R'=K'^\circ$, and let $v'$ be the unique
extension of $v$ to $K'$, so that $v'(\pi)=1$ and
$v'(\pi^{1/m})=1/m$.  The base change $\mathcal U_1\times_RR'$ is
$\Spec R'[x,y]/(xy-\pi)$, whose skeleton still has length
$v'(\pi)=1$.  If instead one \emph{renormalizes} the valuation of $K'$ by
declaring $v''(\pi^{1/m})=1$, then $v''=mv'$ and the same edge has length
$m$.  Edge lengths therefore depend on the normalization of the ground
valuation, not on the field; the ratio of two edge lengths in one skeleton
does not.  This is exactly the ambiguity that disappears over a trivially
valued field only because there the scaling action is treated as part of the
structure (\cref{sec:trivial-valuation-cone}).  Throughout this book the
normalization $v(\pi)=1$ on the base field is in force, and any base change
is accompanied by an explicit statement of which normalization is used.
\end{remark}

To compactify the calculation, consider the projective family
\begin{equation}
\label{eq:running-conic}
             \mathcal C=\{XY=\pi Z^2\}\subset\mathbf P^2_R.
\end{equation}
Its generic fibre is a smooth conic, while its special fibre is the union
$C_0\cup C_1$ of two lines meeting transversely.  The line bundle
$\mathcal L=\mathcal O_{\mathbf P^2_R}(1)|_{\mathcal C}$ has degree $1$ on
each component and degree $2$ on the generic fibre.

\begin{proposition}\label{prop:running-conic-measure}
The model metric on $L=\mathcal L|_{C}$ has Chambert--Loir measure
\begin{equation}
\label{eq:running-conic-measure}
                              \mu_0=\delta_{v_0}+\delta_{v_1}.
\end{equation}
Its mass is $2=\deg_C L$.
\end{proposition}

\begin{proof}
For a semistable curve, the mass at the divisorial point associated to a
component $E$ is the degree of $\mathcal L|_E$.  Both components occur with
multiplicity one and both degrees equal one.  This gives
\eqref{eq:running-conic-measure}; the total degree is preserved in the flat
family.
\end{proof}

Let $\varphi$ be affine on $\Gamma$ with
$\varphi(v_0)=A$, $\varphi(v_1)=B$.  Combining
\eqref{eq:graph-curvature} with \Cref{ex:edge-mass-transport} yields
\begin{equation}
\label{eq:running-perturbed-measure}
 c_1(L,\|\cdot\|_0\rme^{-\varphi})
 =(1+B-A)\delta_{v_0}+(1+A-B)\delta_{v_1}.
\end{equation}
Thus the metric is semipositive precisely when
\begin{equation}
\label{eq:running-nef-interval}
                               |A-B|\leq1.
\end{equation}
Adding a constant to both $A$ and $B$ changes the metric by a constant factor
but leaves its curvature measure unchanged.

\begin{example}[A boundary metric]
Take $B-A=1$.  Formula \eqref{eq:running-perturbed-measure} gives
$2\delta_{v_0}$: all degree has moved to one component.  The metric remains
semipositive, but it lies on the boundary of the nef cone.  If $B-A>1$, the
coefficient at $v_1$ becomes negative, so the corresponding model line bundle
is not nef.
\end{example}

\subsection{The subdivided model, with numbers}
\label{sec:ch08-chain-numbers}

Blowing up the node of $\mathcal C$ inserts the midpoint of the skeleton
(\cref{ex:ch01-node-blowup}); the same computation on the model
$\mathcal C^{(2)}=\{XY=\pi^2Z^2\}$ produces the chain of three lines
$C_0,E,C_2$ over the segment $[0,2]$ analysed in
\cref{ex:chain-nef-functions}.  We record the resulting measures explicitly,
because they are the smallest example in which mass is redistributed among
three atoms.

Let $\cL_0$ be the pullback of $\cO_{\mathbf P^2_R}(1)$, with residual
degrees $1,0,1$ on $C_0,E,C_2$, and let $D=a_0C_0+a_1E+a_2C_2$ be a vertical
$\mathbf Q$-divisor, so that $\varphi_D$ takes the values $a_0,a_1,a_2$ at
the three divisorial points $s=0,1,2$ and is affine on the two unit edges.
The outgoing-slope sums are
\[
 (\Delta\varphi_D)(\{0\})=a_1-a_0,\quad
 (\Delta\varphi_D)(\{1\})=a_0-2a_1+a_2,\quad
 (\Delta\varphi_D)(\{2\})=a_1-a_2 ,
\]
of total mass zero, and therefore
\begin{equation}
\label{eq:ch08-chain-measure}
 c_1\bigl(L,\|\cdot\|_{\cL_0}\rme^{-\varphi_D}\bigr)
 =(1+a_1-a_0)\delta_{s=0}
  +(a_0-2a_1+a_2)\delta_{s=1}
  +(1+a_1-a_2)\delta_{s=2} ,
\end{equation}
of total mass
$(1+a_1-a_0)+(a_0-2a_1+a_2)+(1+a_1-a_2)=2=\deg L$.  These three numbers are
precisely the intersection numbers
$(\cL\cdot C_0)$, $(\cL\cdot E)$, $(\cL\cdot C_2)$ computed in
\cref{ex:chain-nef-functions}, as
\cref{prop:ch08-curve-cl} guarantees.  Three sample values:
\[
\begin{array}{c|c|c}
 (a_0,a_1,a_2) & \text{measure} & \text{semipositive?}\\ \hline
 (0,0,0) & \delta_{s=0}+\delta_{s=2} & \text{yes}\\
 (0,-\tfrac12,0) &
   \tfrac12\delta_{s=0}+\delta_{s=1}+\tfrac12\delta_{s=2} & \text{yes}\\
 (0,1,0) & 2\delta_{s=0}-2\delta_{s=1}+2\delta_{s=2} & \text{no}
\end{array}
\]
In every row the masses sum to $2$.  The middle row is a symmetric
choice putting mass on all three components; more generally
$(0,a_1,0)$ has that property for every $-1<a_1<0$.  The last row shows that a
\emph{concave} kink at an interior vertex destroys semipositivity, which is
the graph form of the requirement that $\varphi_D$ be convex at interior
vertices.

\subsection{Semistable reduction in two examples}
\label{sec:ch08-semistable-reduction}
\index{semistable reduction}

A smooth projective curve over $K$ need not admit a semistable model over
$R$, but it does after a finite base extension; this is the semistable
reduction theorem of Deligne and Mumford, whose higher-dimensional
replacement by alterations is \cite{deJong1996}, and whose abelian-variety
predecessor is Grothendieck's theorem in \cite{SGA7I}.  We compute the two
elliptic cases, which are the model for everything later: the skeleton is a
point precisely when the reduction is potentially good, and a circle
precisely when it is potentially multiplicative.

\begin{example}[Potentially good reduction: a point as skeleton]
\label{ex:ch08-potentially-good}
Assume that the residue characteristic of $K$ is different from $2$ and $3$,
and let
\[
 E:\ y^2=x^3+\pi
\]
be the elliptic curve over $K$ with $v(\pi)=1$.  Its discriminant is
$\Delta=-432\pi^2$, so $v(\Delta)=2$, and its $j$-invariant is $j=0$, so
$v(j)=+\infty\geq0$.  The Weierstrass model $\{y^2=x^3+\pi\}$ over $R$ has special
fibre the cuspidal cubic $y^2=x^3$, which is not nodal: this model is not
semistable, and no semistable model exists over $R$ itself.

Pass to $K'=K(\pi^{1/6})$, and let $v'$ be the extension of $v$, so
$v'(\pi^{1/6})=\tfrac16$.  Substituting
\[
 x=\pi^{1/3}x',\qquad y=\pi^{1/2}y'
\]
turns the equation into $\pi y'^2=\pi x'^3+\pi$, that is,
\[
 y'^2=x'^3+1 ,
\]
whose reduction modulo the maximal ideal of $K'^\circ$ is a smooth plane
cubic, since $\operatorname{char}\widetilde K\neq2,3$.  Hence $E_{K'}$ has
good reduction.  The corresponding model has irreducible smooth special
fibre, its dual complex is a single vertex, and
$\Sk=\{x_{E}\}$ is one divisorial point; the retraction collapses all of
$E_{K'}^{\An}$ to it.  For a line bundle $L$ of degree $d$ with the model
metric, \eqref{eq:ch08-curve-cl-measure} reads
\[
 c_1(L,\|\cdot\|)=d\,\delta_{x_E} ,
\]
one atom of mass $d$.  The exponent $6$ is forced: the base change must make
$v(\Delta)$ divisible by $12$, and $12/\gcd(12,2)=6$.
\end{example}

\begin{example}[Multiplicative reduction: the Tate curve and the circle]
\label{ex:ch08-tate-curve}
\index{Tate curve}
Let $q\in K^\times$ with $\ell:=v(q)>0$ and let
$E_q=\mathbf G_{m,K}^{\An}/q^{\mathbf Z}$ be the Tate curve.  Its
$j$-invariant satisfies $v(j)=-\ell<0$, so the reduction is multiplicative.
The skeleton of $\mathbf G_{m,K}^{\An}$ is the apartment $\mathbf R$ with
coordinate $t=v(z)$, and the subgroup $q^{\mathbf Z}$ acts on it by
translation by $\ell$; hence
\[
 \Sk(E_q)=\mathbf R/\ell\mathbf Z ,
\]
a circle of circumference $\ell$, with the induced retraction
$r(z)=t(z)\bmod\ell$.

If $\ell=n$ is a positive integer, the corresponding semistable model has
special fibre a cycle of $n$ projective lines
$Y_1,\ldots,Y_n$, each meeting the next in one node, and the dual complex is
the circle of circumference $n$ subdivided into $n$ unit edges; this is the
computation quoted at the end of \cref{ex:snc-chain-models}.  Let $L$ have
degree $d$ and let $\cL$ be a model with residual degrees
$d_1,\ldots,d_n$.  Then \eqref{eq:ch08-curve-cl-measure} gives
\[
 c_1(L,\|\cdot\|_{\cL})=\sum_{i=1}^nd_i\,\delta_{v_i},
 \qquad
 \sum_{i=1}^nd_i=d ,
\]
and semipositivity is $d_i\geq0$ for all $i$.

We now solve a Monge--Amp\`ere equation on the circle by hand.  Take
$n=3$ and $d=3$ with the balanced model, $d_1=d_2=d_3=1$, so that
$\mu_0=\delta_{v_1}+\delta_{v_2}+\delta_{v_3}$, and prescribe the target
measure $\mu=3\delta_{v_1}$, of the same mass $3$.  Parametrize the circle by
$s\in\mathbf R/3\mathbf Z$ with $v_1,v_2,v_3$ at $s=0,1,2$.  We must solve
$\Delta\varphi=\mu-\mu_0=2\delta_{s=0}-\delta_{s=1}-\delta_{s=2}$.  By
\eqref{eq:ch08-green-circle} with $\ell=3$, the function
$g_{0,c}$ solving $\Delta g=\delta_0-\delta_c$ is
\[
 g_{0,1}(s)=\begin{cases}\tfrac23s,&0\leq s\leq1,\\[1mm]
                          \tfrac13(3-s),&1\leq s\leq3,\end{cases}
 \qquad
 g_{0,2}(s)=\begin{cases}\tfrac13s,&0\leq s\leq2,\\[1mm]
                          \tfrac23(3-s),&2\leq s\leq3.\end{cases}
\]
Their sum $\varphi=g_{0,1}+g_{0,2}$ is
\[
 \varphi(s)=
 \begin{cases}
  s,&0\leq s\leq1,\\
  1,&1\leq s\leq2,\\
  3-s,&2\leq s\leq3 ,
 \end{cases}
\]
with $\varphi(0)=0$, $\varphi(1)=\varphi(2)=1$.  Its outgoing-slope sums are
$1+1=2$ at $s=0$, $-1+0=-1$ at $s=1$, and $0-1=-1$ at $s=2$; total mass $0$,
as required.  Therefore
\[
 \mu_0+\Delta\varphi
 =(1+2)\delta_{s=0}+(1-1)\delta_{s=1}+(1-1)\delta_{s=2}
 =3\delta_{v_1}=\mu ,
\]
of total mass $3=\deg L$.  The solution $\varphi$ is unique up to an additive
constant by \cref{prop:graph-poisson}, and it is $\mu_0$-psh because $\mu$ is
positive; this is \cref{prop:graph-ma-solution} in a case that can be
checked in three lines.
\end{example}

\begin{remark}[What the two examples separate]
The essential skeleton of \cref{sec:essential-skeleton} is a point in
\cref{ex:ch08-potentially-good} and the whole circle in
\cref{ex:ch08-tate-curve}.  The dichotomy is detected by $v(j)$, and it is
the one-dimensional shadow of the general principle that the essential
skeleton measures the degeneracy of the family, whereas an arbitrary model
skeleton, such as the subdivided segments of
\cref{ex:ch08-An-edge,sec:ch08-chain-numbers}, is a computational device that
can be enlarged at will.
\end{remark}

\subsection{A Green function on the skeleton}

For $q\in(0,1)$, let
\[
 g_q(s)=
 \begin{cases}
  (1-q)s,&0\leq s\leq q,\\
  q(1-s),&q\leq s\leq1.
 \end{cases}
\]
The outgoing slope sum is $-1$ at $q$, $1-q$ at $0$, and $q$ at $1$, so
\begin{equation}
\label{eq:segment-green}
                 \Delta g_q=(1-q)\delta_{v_0}+q\delta_{v_1}-\delta_q.
\end{equation}
This is the potential created by a unit point charge at $q$ balanced by the
endpoint probability measure with barycentric weights.  Formula
\eqref{eq:graph-green} gives its energy:
\[
 -\int_0^1|g_q'|^2\,\rd s=-q(1-q).
\]
The sign is consistent with our Laplacian; the positive Dirichlet energy is
$\int|g_q'|^2=q(1-q)$.

\section{Translation table}
\label{sec:model-translation-table}

The calculations above will reappear in several languages.  This section is
the dictionary, and it is meant to be used in both directions: from an
abstract definition in a later chapter to the computation in this chapter
that fixes its normalization, and from a computation here to the general
theorem it anticipates.  Every entry points to an existing statement.

\subsection{From this chapter to the general theory}
\label{sec:ch08-dict-forward}

\begin{center}
\small
\begin{tabular}{p{.30\textwidth}|p{.26\textwidth}|p{.34\textwidth}}
\textbf{Computed here}&\textbf{Generalized as}&\textbf{Where}\\ \hline
graph Laplacian \eqref{eq:graph-laplacian}
  &$\dprimeddprime$ of a model function on a curve
  &\cref{prop:skeletal-curvature-curve},
   \eqref{eq:cld-curve-graph-laplacian}\\
Green identity \eqref{eq:graph-green}
  &integration by parts and energy pairing
  &\eqref{eq:graph-dirichlet-sign},
   \cref{prop:bj-integration-by-parts}\\
Poisson solution \cref{prop:graph-poisson}
  &Monge--Amp\`ere equation on a curve
  &\cref{prop:graph-ma-solution},
   \eqref{eq:curve-poisson}\\
slope formula \cref{thm:ch08-slope-formula}
  &Poincar\'e--Lelong on a Berkovich curve
  &\cref{chap:cld-poincare-lelong,chap:cld-examples},
   \cref{thm:wanner-curve-comparison}\\
retraction bound \cref{prop:ch08-retraction-max}
  &skeletal characterization of psh functions
  &\cref{prop:skeleton-retraction},
   \cref{thm:bfj-skeletal-characterization}\\
curve model measure \eqref{eq:ch08-curve-cl-measure}
  &Chambert--Loir measure of a model metric
  &\cref{sec:metrics-cl-measure},
   \cref{thm:cl-model-properties}\\
model function of a vertical divisor
  &PL functions and vertical $\mathbf Q$-Cartier divisors
  &\cref{prop:skeleton-model-pa},
   \cref{thm:skeleton-pl-realization}\\
two-legged-tree measure \eqref{eq:ch08-p1-tripod-measure}
  &atomic measures on a dual complex
  &\cref{chap:bfj-model-forms},
   \eqref{eq:ch20-graph-laplacian}\\
trivially valued tree
   (\cref{prop:trivial-p1-tree})
  &the analytification $X^{\An}$ over a trivially valued field
  &\cref{def:trivial-analytification},
   \cref{ex:triv-p1-tree}\\
scaling $v\mapsto tv$
  &the cone structure and normalized links
  &\eqref{eq:valuation-scaling},
   \cref{def:normalized-link}\\
$\varphi_{\mathfrak a}=-v(\mathfrak a)$
   (\cref{prop:ch08-ideal-potential})
  &homogeneous functions and b-divisors
  &\cref{def:homogeneous-function},
   \cref{prop:cartier-homogeneous-pl}\\
Gauss extension \cref{ex:ch08-gauss-ray}
  &the section $\sigma:X^{\An}\to X_K^{\An}$
  &\eqref{eq:gauss-extension-triv},
   \eqref{eq:triv-gauss-section}\\
$\MA_{\mathbf R}$ of a polyhedral $F$
   (\cref{lem:ch08-alexandrov-polyhedral})
  &toric Monge--Amp\`ere comparison
  &\cref{thm:toric-bfj-alexandrov},
   \eqref{eq:toric-bfj-alexandrov}\\
toric convexity \cref{sec:toric-convex}
  &invariant psh functions are convex
  &\cref{thm:bj-toric-psh-convex},
   \eqref{eq:bj-toric-ma}\\
circle skeleton \cref{ex:ch08-tate-curve}
  &essential skeleta and hybrid limits
  &\cref{sec:essential-skeleton},
   \cref{ex:hybrid-nodal-curve}
\end{tabular}
\end{center}

\subsection{From the general theory back to this chapter}
\label{sec:ch08-dict-backward}

The reverse dictionary is the one to consult when a later definition looks
unmotivated or a sign looks arbitrary.

\begin{center}
\small
\begin{tabular}{p{.30\textwidth}|p{.30\textwidth}|p{.30\textwidth}}
\textbf{Later object}&\textbf{Smallest instance}&\textbf{Computed in}\\ \hline
$\theta$-psh model function
  &$\varphi$ affine on $[0,1]$ with $|A-B|\leq1$
  &\eqref{eq:running-nef-interval},
   \cref{ex:edge-mass-transport}\\
Chambert--Loir measure $c_1(L,\|\cdot\|)^n$
  &$\delta_{\zeta_{0,1}}$ on $\mathbf P^1$
  &\eqref{eq:p1-standard-measure},
   \cref{ex:ch08-toric-pn}\\
mass identity $\int\MA=(L^n)$
  &$1+1=2$ on the running conic
  &\cref{prop:running-conic-measure},
   \cref{ex:degree-conservation-conic}\\
regularization by model metrics
  &$\varepsilon\log(1+\rme^{-du/\varepsilon})\to\max\{0,-du\}$
  &\cref{sec:toric-convex},
   \cref{ex:ch01-equilibrium-metric}\\
Fubini--Study function
  &$\max\{0,\log|T|-\lambda\}$
  &\cref{ex:ch08-p1-fs},
   \cref{sec:metrics-fs}\\
test configuration and its measure
  &$\delta_{c\,\ord_0}$ on the trivially valued tree
  &\cref{ex:ch08-trivial-one-ray},
   \cref{ex:ch01-product-tc}\\
homogeneous potential of an ideal
  &$\max\{-2w_1,-3w_2\}$ on the quadrant
  &\cref{sec:trivial-valuation-cone},
   \cref{ex:ch08-monomial-multiplicity}\\
support function and polytope
  &$\check h_{[0,d]}(u)=\max\{0,-du\}$
  &\cref{sec:toric-convex},
   \eqref{eq:appJ-toric-recession}\\
factor $n!$ between $\MA_{\mathbf R}$ and $\MA$
  &$2!\cdot\tfrac32=3=(2H-E)^2$
  &\cref{ex:ch08-toric-hirzebruch},
   \eqref{eq:appJ-toric-degree}\\
semistable model and its dual complex
  &cycle of $n$ lines, circle of circumference $n$
  &\cref{ex:ch08-tate-curve},
   \cref{ex:snc-chain-models}\\
model comparison of two theories
  &one atom computed by slopes and by degrees
  &\cref{ex:ch08-p1-interval},
   \cref{prop:model-core-preview}
\end{tabular}
\end{center}

\subsection{The three running examples across the dictionary}
\label{sec:ch08-dict-running}

\begin{itemize}
\item \emph{The nodal degeneration} of \cref{sec:ch01-running-nodal} appears
here as the annulus and its interval skeleton
(\cref{prop:ch08-annulus-skeleton}), as the conic $\mathcal C$ with measure
$\delta_{v_0}+\delta_{v_1}$ (\cref{prop:running-conic-measure}), as the
subdivided chain with the three-atom measure \eqref{eq:ch08-chain-measure},
and as the Tate curve with a circle skeleton
(\cref{ex:ch08-tate-curve}).  It returns in
\cref{chap:cld-examples,chap:bfj-examples}, in
\cref{ex:cma-semistable-curve}, and in \cref{ex:hybrid-nodal-curve}.

\item \emph{The toric line} of \cref{sec:ch01-running-toric} appears here as
the apartment of \cref{sec:berkovich-p1}, as
$\check h_{[0,d]}=\max\{0,-du\}$, and as the one-dimensional case of
\cref{thm:toric-ma-comparison}; the two-break computation
\cref{ex:ch01-toric-two-breaks} is the case $n=1$ of
\cref{lem:ch08-alexandrov-polyhedral}.  Its higher-dimensional companions are
\cref{ex:ch08-toric-p1p1,ex:ch08-toric-hirzebruch,ex:ch08-toric-weighted},
and it returns in \cref{thm:toric-bfj-alexandrov,thm:bj-toric-psh-convex} and
in \cref{ex:verified-toric-segment}.

\item \emph{The product test configuration} of \cref{sec:ch01-running-tc}
appears here as \cref{ex:ch08-trivial-one-ray} with $c=1$, and its
homogeneous potential as \cref{ex:ch08-truncation}.  It returns in
\cref{chap:bj-test-configurations,chap:rays-filtrations} and in
\cref{ex:dict-product-action}.
\end{itemize}

\subsection{The mass ledger}
\label{sec:ch08-mass-ledger}

Every Monge--Amp\`ere computation of this chapter has had its total mass
checked against a degree.  The checks are collected here, so that a reader
who has reproduced them knows that all conventions are consistent.

\begin{center}
\small
\begin{tabular}{p{.40\textwidth}|p{.24\textwidth}|p{.24\textwidth}}
\textbf{Computation}&\textbf{Mass}&\textbf{Degree}\\ \hline
$\cO(1)$ on $\mathbf P^1$, standard metric,
  \eqref{eq:p1-standard-measure}
  &$1$&$\deg\cO(1)=1$\\
$\cO(1)$ on $\mathbf P^1$, one blowup,
  \eqref{eq:ch08-p1-interval-measure}
  &$(1+a)+(-a)=1$&$1$\\
$\cO(1)$ on $\mathbf P^1$, two-legged tree,
  \eqref{eq:ch08-p1-tripod-measure}
  &$(1+a_0+a_1)-a_0-a_1=1$&$1$\\
$\cO(2)$ on $\mathbf P^1$, two atoms,
  \cref{ex:ch08-p1-degree-two}
  &$1+1=2$&$2$\\
running conic, \eqref{eq:running-conic-measure}
  &$1+1=2$&$\deg_C L=2$\\
running conic, perturbed,
  \eqref{eq:running-perturbed-measure}
  &$(1+B-A)+(1+A-B)=2$&$2$\\
chain model, \eqref{eq:ch08-chain-measure}
  &$2$ for all $(a_0,a_1,a_2)$&$2$\\
Tate curve, $n=3$, \cref{ex:ch08-tate-curve}
  &$3$&$\deg L=3$\\
$\mathbf P^n$, \cref{ex:ch08-toric-pn}
  &$n!\cdot\tfrac1{n!}=1$&$(\cO(1)^n)=1$\\
$\mathbf P^1\times\mathbf P^1$, $\cO(a,b)$,
  \cref{ex:ch08-toric-p1p1}
  &$2!\,ab=2ab$&$(aH_1+bH_2)^2=2ab$\\
$\mathbf P^1\times\mathbf P^1$, $F_\lambda$,
  \cref{ex:ch08-toric-p1p1}
  &$2(\tfrac12+\tfrac12)=2$&$2$\\
Hirzebruch $\mathbf F_1$, \cref{ex:ch08-toric-hirzebruch}
  &$2(1+\tfrac12)=3$&$(2H-E)^2=3$\\
$\mathbf P(1,1,2)$, $\cO(2)$,
  \cref{ex:ch08-toric-weighted}
  &$2!\cdot1=2$&$(\cO(2)^2)=2$\\
$\mathbf P(1,1,2)$, $\cO(1)$,
  \cref{ex:ch08-toric-weighted}
  &$2!\cdot\tfrac14=\tfrac12$&$(\cO(1)^2)=\tfrac12$\\
trivially valued $\mathbf P^1$, one ray,
  \cref{ex:ch08-trivial-one-ray}
  &$1$&$\deg\cO(1)=1$\\
trivially valued $\mathbf P^1$, two rays,
  \cref{ex:ch08-trivial-two-rays}
  &$\tfrac12+\tfrac12=1$&$1$\\
monomial ideal $(x^2,y^3)$,
  \cref{ex:ch08-monomial-multiplicity}
  &$2!\cdot3=6$&$e(\mathfrak a)=6$
\end{tabular}
\end{center}

Two further checks are of a different type: the slope formula was verified
atom by atom in \cref{ex:ch08-newton-cubic}, where both sides have total mass
zero, and the Laplacian of every model function has total mass zero by
\eqref{eq:lap-mass-zero}, which is what makes the mass of a curvature measure
depend only on the line bundle.

\subsection{The three-column summary}
\label{sec:ch08-dict-three-column}

Finally, the compact form of the dictionary.
\begin{center}
\small
\begin{tabular}{p{.27\textwidth}|p{.31\textwidth}|p{.29\textwidth}}
\textbf{Model or analytic object}&\textbf{Skeleton or tropical object}
  &\textbf{Curvature object}\\ \hline
vertical divisor on a curve model
  &integral PL function on $\Gamma$
  &$\Delta_\Gamma\varphi$\\
nef model metric on a curve
  &$\mu_0+\Delta_\Gamma\varphi\geq0$
  &positive atomic measure\\
rational function
  &$-\log|f|$ with integral slopes
  &$-\delta_{\operatorname{div}(f)}$\\
toric semipositive metric
  &convex $F\sim\check h_P$
  &$n!\operatorname{MA}_{\mathbf R}(F)$\\
toric model metric
  &rational PL convex $F$
  &atoms of mass $n!\operatorname{vol}(\partial F)$\\
trivially valued model ideal
  &$-v(\mathfrak a)$ on the valuation cone
  &homogeneous PL data
\end{tabular}
\end{center}

Four normalization checks should be remembered.
\begin{enumerate}
\item $v=-\log|\cdot|$, so $v(f)$ and $\log|f|$ have opposite Laplacians.
\item The graph Laplacian is the sum of outgoing slopes.
\item Alexandrov measure has total mass $\operatorname{vol}(P)$, whereas the
  geometric top wedge has mass $n!\operatorname{vol}(P)=(L^n)$.
\item No probability normalization is applied unless it is written
  explicitly.
\end{enumerate}

\section{Source note}

Potential theory on metric graphs, in the form of
\cref{sec:metric-graphs}, is classical; the version used here---outgoing
slopes, Green's identity, the Poisson problem, and the maximum
principle---is the one that appears in tropical geometry, for which
\cite{MR18Tropical} is a convenient systematic reference, and in the
potential theory of Berkovich curves developed by Thuillier \cite{Thu05}
and, for $\mathbf P^1$, by Baker and Rumely, whose monograph is named rather
than cited by key because no proof here takes an input from it.

The primary source for \cref{sec:metric-graphs} is Thuillier's
\cite[\S1.2]{Thu05}, and we have quoted his numbering throughout: the
Laplacian is his \cite[\S1.2.1]{Thu05}, the maximum principle his
\cite[Prop.~1.2.9, Cor.~1.2.10]{Thu05}, Green's identity and the vanishing of
the total mass his \cite[Prop.~1.2.13, Cor.~1.2.14]{Thu05}, the Dirichlet
problem his \cite[Prop.~1.2.15, Rem.~1.2.16]{Thu05}, the Poisson problem his
\cite[Prop.~1.2.17, Lem.~1.2.18]{Thu05}, and the invariance under retraction
his \cite[Prop.~1.2.20]{Thu05}.  His \cite[\S1.2.5]{Thu05} extends $\ddc$ to
distributions on a polyhedron and identifies the extension with Zhang's
operator on the space $\mathrm F(S)$ of edgewise smooth functions
\cite[Prop.~1.2.21]{Thu05}, up to the sign $\ddc=-\Delta$ that his
comparison records; that is the only sign conversion needed anywhere in the
comparison, and it does not affect us because we never use Zhang's
normalization.  What \cite{Thu05} adds beyond
\cref{sec:metric-graphs}---the sheaves $\mathcal H_C$ and
$\mathcal{SH}_C$ on the whole of $C^{\An}$, Perron's method, currents, and
the approximation theorem---is used in \cref{chap:cld-examples}; see
\cref{rem:ch08-harmonic-vs-log,rem:ch08-thuillier-normalization} for the two
points at which his theory is sharper than, or normalized differently from,
what is proved here.
\Cref{thm:ch08-slope-formula} is the elementary model-theoretic form of the
slope formula; the analytic statement valid on all of $C^{\An}$, with the
type-$\mathrm{IV}$ points and the branches off the skeleton, belongs to that
literature and is compared with the Chambert--Loir--Ducros calculus in
\cite{Wanner19} and, in the recent curve-test framework, in
\cite{GR25PL,GR25Classical}; a Poincar\'e lemma for the associated real forms
is \cite{Jell16}.  The convexity facts used in \cref{sec:toric-convex} are
standard and may be found in \cite{Roc70,Sch14}; toric background is
\cite{Fultor,CLS11}, and the arithmetic toric formalism closest to our
normalization is \cite{BGPS14}.  Semistable reduction for curves is the
theorem of Deligne and Mumford, resting on Grothendieck's semistable
reduction for abelian varieties in \cite{SGA7I}; the alteration replacement
in higher dimension is \cite{deJong1996}.  The Tate curve computation of
\cref{ex:ch08-tate-curve} is the analytic uniformization of an elliptic curve
with multiplicative reduction; a formal account in the language used here is
\cite{BoschLutkebohmert1993}.  Intersection-theoretic inputs are quoted from
\cite{Ful} through \cref{chap:models}.

The graph calculations are the model core of potential theory on Berkovich
curves; see also \cite{Berk12}.  Historically the order was the reverse of
the order of exposition here: Rumely's capacity theory on algebraic curves
and Kani's potential-theoretic reading of intersection theory on an
arithmetic surface came first, Thuillier \cite{Thu05} then rebuilt both on
Berkovich spaces by reducing everything to \cref{sec:metric-graphs} through
Bosch--L\"utkebohmert semistable reduction \cite[Thm.~2.3.8]{Thu05}, and
Baker and Rumely afterwards developed the projective line in the same spirit.
Thuillier's own comparison with these three approaches, and with the
valuative-tree potential theory of Favre and Jonsson, is
\cite[\S\S5.1--5.4]{Thu05}; it is summarized in the historical note of
\cref{chap:cld-examples}.  The equality between the curve-theoretic and
CLD notions of
subharmonicity will be discussed in \cref{chap:cld-examples}.  The atomic model measure is
the construction of Chambert--Loir \cite{CL06}; its recovery from CLD Chern
currents is treated later using \cite{CLD12}.  The global model/psh framework
for nontrivially valued fields is developed in \cite{BFJ16}.  The toric
real/non-Archimedean comparison is represented at three distinct levels:
the model measure of \cite{CL06}, the calibrated polystable-face theorem of
\cite{Vil20}, and the global invariant tropical theory of \cite{BGJKK25}.
The delta-form and tropical-cycle refinements of \cite{GK17,GK19}, and the
harmonic-tropical enlargement of the form sheaf in \cite{GJR26}, will be
introduced in \cref{chap:delta-forms} rather than folded silently into the CLD calculus.
The curve-test theory of \cite{GR25PL,GR25Classical} enters in
\Cref{chap:cld-poincare-lelong,chap:cld-examples}, where it will be
compared---not identified by
definition---with CLD psh functions.  The trivially valued cone and its
homogeneous psh functions are developed systematically in \cite{BJ22}.
